\documentclass[a4paper, 12pt, twoside]{article}
\usepackage[utf8]{inputenc}
\usepackage[T1]{fontenc}
\usepackage[french]{babel}
\usepackage{lmodern}
\usepackage[top=2.5cm, bottom=2cm,
            left=3cm, right=2.5cm,
            headheight=15pt]{geometry}
\usepackage{graphicx}
\usepackage{eso-pic}
\usepackage{array}
\usepackage{hyperref}
\usepackage{listings}
\usepackage{xcolor}
\usepackage{amsmath}
\usepackage{amssymb}
\usepackage{booktabs}
\usepackage{microtype}
\setlength{\emergencystretch}{3em}
\input{pagedegarde}

\lstset{
  basicstyle=\footnotesize\ttfamily,
  keywordstyle=\color{blue},
  commentstyle=\color{gray},
  breaklines=true,
  frame=single,
  numbers=none,
}

\title{Op\'{e}ration No\"{e}l}

\datedebut{janvier 2025}
\datefin{mai 2025}

\membrea{Bekkaripro Salah}
\membreb{}

\begin{document}
\pagedegarde

\section*{Remerciements}
Je remercie les enseignants du d\'{e}partement MIAGE de l'Universit\'{e} Paris Nanterre pour
l'encadrement de ce projet, ainsi que les \'{e}quipes OpenStreetMap, Open-Meteo et ADEME pour la
mise \`{a} disposition gratuite de leurs donn\'{e}es et APIs.

\newpage
\tableofcontents
\newpage

%=======================================================
\section{Introduction}
%=======================================================

Ce rapport pr\'{e}sente le projet \textbf{Op\'{e}ration No\"{e}l}, une plateforme web d'optimisation
logistique urbaine d\'{e}velopp\'{e}e dans le cadre du cours Graphes et Open Data de la Licence MIAGE.

\subsection{Probl\'{e}matique}

La probl\'{e}matique retenue est la suivante~:
\begin{quote}
\emph{Comment exploiter les graphes et l'open data pour organiser des tourn\'{e}es de livraison
du P\`{e}re No\"{e}l r\'{e}alistes, optimis\'{e}es et adaptables aux contraintes du terrain~?}
\end{quote}

Cette question m\^{e}le deux dimensions du cours~: les \textbf{graphes}, car le r\'{e}seau routier
est mod\'{e}lis\'{e} sous forme de n\oe uds et d'ar\^{e}tes pond\'{e}es, et l'\textbf{open data},
car les tourn\'{e}es sont construites \`{a} partir de donn\'{e}es r\'{e}elles comme OpenStreetMap,
la m\'{e}t\'{e}o ou le relief. Le probl\`{e}me sous-jacent est un VRPTW (Vehicle Routing Problem
with Time Windows), NP-difficile~: le nombre de solutions possibles cro\^{i}t tr\`{e}s vite avec
le nombre de clients, ce qui rend une recherche exhaustive impraticable.

\subsection{Approche retenue}

Le projet combine trois couches~:
\begin{enumerate}
  \item Des donn\'{e}es \textbf{Open Data} issues de plusieurs sources publiques
        (OpenStreetMap, Open-Meteo, Overpass, ADEME). OSRM n'est conserv\'{e}
        que comme outil historique pour le jeu de r\'{e}f\'{e}rence.
  \item Un \textbf{solveur de recherche op\'{e}rationnelle} (Google OR-Tools) configurable par profil
        de strat\'{e}gie.
  \item Des \textbf{algorithmes de graphes} impl\'{e}ment\'{e}s de bout en bout (Dijkstra, A*,
        Floyd-Warshall, K-Means, Louvain, centralit\'{e}), expos\'{e}s \`{a} la fois comme outils de
        calcul et comme supports p\'{e}dagogiques visualisables.
\end{enumerate}

%=======================================================
\section{Environnement de travail}
%=======================================================

\subsection{Stack technique}

\begin{table}[h]
  \begin{center}
    \begin{tabular}{|l|l|}
      \hline
      \textbf{Composant} & \textbf{Technologie} \\
      \hline
      Syst\`{e}me & Linux Ubuntu 22.04 LTS \\
      Backend & Python 3.10, FastAPI, Uvicorn \\
      Solveur & Google OR-Tools (\texttt{ortools}) \\
      Graphes & OSMnx $\geq$2.0.0, NetworkX \\
      Donn\'{e}es & NumPy, Pandas \\
      Frontend & Next.js 14, TypeScript, React 18 \\
      Cartographie & Leaflet + React-Leaflet, Mapbox GL \\
      Base de donn\'{e}es & SQLite (module standard Python) \\
      D\'{e}ploiement & Docker Compose \\
      Tests & Pytest (20 fichiers) + Playwright (E2E) \\
      Versionnement & Git \\
      \hline
    \end{tabular}
  \end{center}
  \caption{Stack technique du projet}
  \label{tab:stack}
\end{table}

\subsection{Architecture g\'{e}n\'{e}rale}

Le projet est s\'{e}par\'{e} en deux services Docker~:
\begin{itemize}
  \item \textbf{Backend} (port 8000)~: FastAPI expose 70 routes REST et 2 WebSockets.
        Toute la logique m\'{e}tier (g\'{e}n\'{e}ration de zone, calcul de matrices, solveur VRP,
        algorithmes de graphes, mod\`{e}le d'apprentissage) est concentr\'{e}e dans
        \texttt{backend/app/services.py} et les scripts Python.
  \item \textbf{Frontend} (port 3000)~: Next.js 14 sert l'interface web. La cartographie
        est rendue c\^{o}t\'{e} client avec Leaflet. Les requ\^{e}tes API utilisent TanStack Query
        pour la mise en cache et la synchronisation.
\end{itemize}

\subsection{M\'{e}thodologie de d\'{e}veloppement assist\'{e} par IA}

Le d\'{e}veloppement a \'{e}t\'{e} assist\'{e} par plusieurs agents IA~: Claude Code pour
l'organisation et la coh\'{e}rence d'ensemble, Gemini CLI pour les parties math\'{e}matiques
(score, matrices, heuristiques) et Codex pour l'impl\'{e}mentation, les tests et les
corrections. Ces outils ont acc\'{e}l\'{e}r\'{e} le travail, mais les choix finaux, les tests
et l'int\'{e}gration sont rest\'{e}s sous contr\^{o}le du d\'{e}veloppeur.

\subsection{Organisation des fichiers}

\begin{table}[h]
  \begin{center}
    \begin{tabular}{|l|l|}
      \hline
      \textbf{R\'{e}pertoire} & \textbf{Contenu} \\
      \hline
      \texttt{backend/app/} & API REST, logique m\'{e}tier, persistance SQLite \\
      \texttt{scripts/} & G\'{e}n\'{e}ration OSM, algorithmes de graphes, benchmark \\
      \texttt{final\_scripts/} & Solveur VRPTW principal \\
      \texttt{core\_data/} & Graphe de r\'{e}f\'{e}rence (paris5), matrices, donn\'{e}es clients \\
      \texttt{frontend/app/} & Pages Next.js (solver, mission, versus, leaderboard) \\
      \texttt{tests/} & Suite pytest + Playwright \\
      \hline
    \end{tabular}
  \end{center}
  \caption{Organisation du d\'{e}p\^{o}t}
  \label{tab:structure}
\end{table}

\subsection{Lien du d\'{e}p\^{o}t}

Le code source complet du projet est versionn\'{e} sur GitHub~:
\url{https://github.com/salahbekkaripro/Santa_Delivery}.

%=======================================================
\section{Description du projet et objectifs}
%=======================================================

\subsection{Objectif}

Op\'{e}ration No\"{e}l est une application web ludique o\`{u} l'utilisateur joue le r\^{o}le d'un
logisticien du P\`{e}re No\"{e}l. Il g\'{e}n\`{e}re une mission de livraison dans n'importe quelle
ville du monde, peut tenter de construire manuellement ses tourn\'{e}es, puis compare ses d\'{e}cisions
\`{a} celles d'un solveur d'optimisation. Un score et un rang r\'{e}compensent l'efficacit\'{e}
de la solution produite.

\subsection{P\'{e}rim\`{e}tre g\'{e}ographique}

L'application est mondiale~: l'utilisateur peut saisir n'importe quelle adresse (Paris, Tokyo,
Montr\'{e}al, etc.). Un jeu de donn\'{e}es de r\'{e}f\'{e}rence pr\'{e}-calcul\'{e}, nomm\'{e}
\texttt{paris5} dans le d\'{e}p\^{o}t, couvre une zone parisienne de test avec un d\'{e}p\^{o}t
central et 10 clients.

\subsection{Fonctionnalit\'{e}s impl\'{e}ment\'{e}es}

\begin{enumerate}
  \item G\'{e}n\'{e}ration d'une mission~: t\'{e}l\'{e}chargement OSM, s\'{e}lection de clients,
        calcul des matrices de co\^{u}t.
  \item R\'{e}solution automatique par solveur VRPTW (3 profils disponibles dans l'interface,
        7 dans le backend).
  \item Mode de jeu humain~: l'utilisateur choisit ses it\'{e}raires segment par segment.
  \item Incidents dynamiques~: routes bloqu\'{e}es, re-planification automatique.
  \item Analyse structurelle du graphe~: centralit\'{e}, robustesse, secteurs de livraison.
  \item Visualisation p\'{e}dagogique~: Dijkstra, A* bidirectionnel, Floyd-Warshall pas-\`{a}-pas.
  \item Score et classement~: 5 rangs (S, A, B, C, D).
  \item Mode versus exp\'{e}rimental~: duel en temps r\'{e}el entre deux joueurs (WebSockets).
  \item Recommandation contextuelle~: suggestion du profil IA le plus adapt\'{e} selon
        l'historique disponible.
\end{enumerate}

\subsection{Choix de p\'{e}rim\`{e}tre}

Au d\'{e}but du projet, l'objectif \'{e}tait de d\'{e}velopper en parall\`{e}le de nombreuses
fonctionnalit\'{e}s~: solveur, mode humain, versus, social, campagne, classement et apprentissage.
La r\'{e}alisation a montr\'{e} que le c\oe ur du sujet Graphes et Open Data se trouvait surtout
dans la g\'{e}n\'{e}ration des graphes, le calcul des matrices, les contraintes de tourn\'{e}es
et l'optimisation. Les modules sociaux, versus et campagne ont donc \'{e}t\'{e} conserv\'{e}s
comme extensions secondaires~: ils existent dans l'application, mais ils ne sont pas consid\'{e}r\'{e}s
comme le c\oe ur finalis\'{e} du projet.

%=======================================================
\section{Parcours utilisateur sur le site}
%=======================================================

Cette section d\'{e}crit le parcours concret sur l'interface, en liant chaque action \`{a} la
cha\^{i}ne technique qu'elle d\'{e}clenche.

\subsection{\'{E}tape 1 -- Saisie de l'adresse}

L'utilisateur acc\`{e}de \`{a} la page \textbf{Solveur de tourn\'{e}es}. Il tape librement une adresse
dans un champ de texte avec auto-compl\'{e}tion~: chaque frappe d\'{e}clenche (apr\`{e}s un
d\'{e}lai de 280\,ms) un appel \`{a} l'\textbf{API Mapbox Geocoding}
(\texttt{https://api.mapbox.com/geocoding/v5/mapbox.places/}), qui renvoie
des suggestions de lieux (adresses, quartiers, POI).

Lorsqu'une suggestion est s\'{e}lectionn\'{e}e, les coordonn\'{e}es GPS (latitude, longitude) du
point central sont v\'{e}rouill\'{e}es. Une mini-carte Leaflet affiche imm\'{e}diatement le cercle
de la zone de recherche, dont le rayon est r\'{e}glable de 0,5 \`{a} 30\,km.

\subsection{\'{E}tape 2 -- Configuration de la mission}

L'utilisateur param\`{e}tre~:
\begin{itemize}
  \item Le nombre de colis \`{a} livrer (8 \`{a} 2\,000).
  \item Le budget total et le co\^{u}t par tra\^{i}neau.
  \item Le \textbf{profil IA}~: trois choix sont pr\'{e}sent\'{e}s dans l'interface~:
    \textbf{Express} (vitesse maximale, bonus +2\,pts), \textbf{\'{E}colo} (distance minimale,
    bonus +3\,pts), \textbf{Prudent} (marges de s\'{e}curit\'{e}, bonus +4\,pts). Ces profils
    correspondent \`{a} trois des sept pr\'{e}sets d\'{e}finis dans le backend~; les autres
    (\textit{Agressive}, \textit{Opportuniste}, \textit{Championne}, \textit{Championne Secteurs})
    ne sont pas expos\'{e}s dans l'UI actuelle.
  \item Des options avanc\'{e}es~: incidents al\'{e}atoires, relief SRTM, CO2 ADEME, limite
        de tra\^{i}neaux.
\end{itemize}

\subsection{\'{E}tape 3 -- G\'{e}n\'{e}ration de la mission (pipeline c\^{o}t\'{e} serveur)}

Quand l'utilisateur lance le calcul de la tourn\'{e}e, le frontend
envoie deux requ\^{e}tes successives au backend~:

\paragraph{Requ\^{e}te 1 --- Cr\'{e}ation de mission.}
Le backend ex\'{e}cute le pipeline suivant (\texttt{scripts/generator\_engine.py})~:

\begin{enumerate}
  \item \textbf{T\'{e}l\'{e}chargement OSM}~: OSMnx appelle les APIs Nominatim et Overpass pour
        t\'{e}l\'{e}charger le r\'{e}seau routier de la zone d\'{e}finie par le point central et
        le rayon. Le graphe est un MultiDiGraph dirig\'{e} repr\'{e}sentant les intersections
        (n\oe uds) et les tron\c{c}ons de voie (ar\^{e}tes).
  \item \textbf{Annotation des ar\^{e}tes}~: chaque ar\^{e}te re\c{c}oit quatre attributs calcul\'{e}s
        localement \`{a} partir de la g\'{e}om\'{e}trie OSM~: temps de trajet, distance, \'{e}missions
        CO2 (ou facteur ADEME si l'option est activ\'{e}e), et score de risque selon le type de voie.
  \item \textbf{S\'{e}lection des points}~: un d\'{e}p\^{o}t central et $n$ clients sont s\'{e}lectionn\'{e}s
        parmi les n\oe uds du graphe. Les noms des clients sont enrichis avec de vrais noms
        d'\'{e}tablissements r\'{e}cup\'{e}r\'{e}s via l'API Overpass (commerces, \'{e}quipements).
  \item \textbf{Calcul des matrices}~: pour chaque paire de points, NetworkX calcule le plus court
        chemin sur le graphe OSM annot\'{e}. Cela produit 5 matrices $n \times n$~: temps, distance,
        CO2, risque, et une matrice composite pond\'{e}r\'{e}e.
  \item \textbf{M\'{e}t\'{e}o}~: la condition m\'{e}t\'{e}orologique choisie (ou r\'{e}elle via
        Open-Meteo) est appliqu\'{e}e comme multiplicateur sur la matrice de temps.
  \item \textbf{Congestion horaire}~: si une heure de d\'{e}part est sp\'{e}cifi\'{e}e, un facteur
        de congestion (ex. 1,8$\times$ \`{a} 18h) est appliqu\'{e}.
\end{enumerate}

L'interface affiche en temps r\'{e}el les \'{e}tapes de chargement (\og Chargement du graphe OSM \fg{},
\og G\'{e}n\'{e}ration des points \fg{}, \og Calcul de la matrice \fg{}, etc.).

\paragraph{Requ\^{e}te 2 --- R\'{e}solution VRPTW.}
Avec les matrices calcul\'{e}es, le backend appelle OR-Tools pour r\'{e}soudre le VRPTW. La
strat\'{e}gie de r\'{e}solution (premi\`{e}re solution, m\'{e}taheuristique, limite de temps) est
s\'{e}lectionn\'{e}e automatiquement selon le profil IA choisi.

\subsection{\'{E}tape 4 -- Affichage des r\'{e}sultats}

La carte Leaflet affiche les tourn\'{e}es optimis\'{e}es, color\'{e}es par tra\^{i}neau. L'utilisateur
voit~:
\begin{itemize}
  \item Le \textbf{score} (0--100) avec ses composantes (temps, CO2, budget, couverture) et
        le rang obtenu.
  \item Les KPIs~: temps total, distance, CO2 \'{e}conom\'{e}, nombre de tra\^{i}neaux utilis\'{e}s,
        gain par rapport \`{a} la solution na\"ive.
  \item Le \textbf{d\'{e}tail des tourn\'{e}es}~: ordre des clients par tra\^{i}neau, avec ETA
        (heure d'arriv\'{e}e estim\'{e}e) et poids charg\'{e}.
  \item La \textbf{recherche automatique du nombre de tra\^{i}neaux}~: le solveur teste
        plusieurs valeurs de $k$ (nombre de v\'{e}hicules) et retient celle qui maximise
        le meilleur compromis entre co\^{u}t op\'{e}rationnel et clients non livr\'{e}s.
\end{itemize}

\subsection{\'{E}tape 5 -- Options avanc\'{e}es}

Depuis la page de r\'{e}sultats, l'utilisateur peut~:
\begin{itemize}
  \item \textbf{Simuler un incident}~: une route al\'{e}atoire est bloqu\'{e}e et le solveur
        re-planifie en temps r\'{e}el.
  \item \textbf{Acc\'{e}der aux m\'{e}triques du graphe}~: densit\'{e}, composantes fortement
        connexes, top-5 des n\oe uds critiques par centralit\'{e} d'interm\'{e}diarit\'{e}.
  \item \textbf{Consulter le d\'{e}bogueur de strat\'{e}gie}~: profil utilis\'{e}, param\`{e}tres
        OR-Tools, statut du solveur.
\end{itemize}

%=======================================================
\section{Architecture fonctionnelle du site}
%=======================================================

Au-del\`{a} du solveur, le projet est organis\'{e} comme une plateforme compl\`{e}te. Chaque page
du site correspond \`{a} un usage distinct, mais toutes reposent sur les m\^{e}mes briques~:
missions persist\'{e}es, graphes OSM, matrices de co\^{u}t, scores et comptes joueurs.

Les modules principaux sont~: \textbf{Solveur} (configuration et lancement), \textbf{Mission}
(carte, r\'{e}sultats, d\'{e}briefing), \textbf{Mode humain} (construction manuelle de tourn\'{e}es),
\textbf{Explore} (visualisation Dijkstra, A*, Floyd-Warshall, 2-opt, OR-opt, ILS),
\textbf{Data} (sources open data et caches) et \textbf{Leaderboard} (scores solo).

Les modules \textbf{Versus}, \textbf{Social}, \textbf{Authentification} et \textbf{Campagne}
compl\`{e}tent la plateforme, mais ils sont secondaires par rapport au sujet principal.
Le Versus et la Campagne sont donc pr\'{e}sent\'{e}s comme des extensions exp\'{e}rimentales~:
ils rendent le site plus vivant, mais le noyau finalis\'{e} reste la partie graphe,
open data et solveur.

%=======================================================
\section{Sources Open Data et services externes utilis\'{e}s}
%=======================================================

\subsection{Connexion entre les APIs}

Le sch\'{e}ma ci-dessous (reproduit de la page \textit{Data} du site) illustre le flux de donn\'{e}es
entre les APIs externes et les modules internes~:

\begin{center}
\begin{tabular}{lcl}
  \textbf{Source ou service} & $\rightarrow$ & \textbf{Donn\'{e}e produite} \\
  \hline
  Mapbox Geocoding & $\rightarrow$ & Coordonn\'{e}es GPS du point central \\
  OpenStreetMap (OSMnx) & $\rightarrow$ & Graphe routier (n\oe uds + ar\^{e}tes) \\
  Overpass API & $\rightarrow$ & Noms r\'{e}els des commerces (clients) \\
  Open-Meteo & $\rightarrow$ & Facteur m\'{e}t\'{e}o $\times$ matrice de temps \\
  ADEME ImpactCO2 & $\rightarrow$ & Facteur CO2 en g/km (optionnel) \\
  \textit{(calcul local)} & $\rightarrow$ & Matrices $n \times n$ via NetworkX \\
  \hline
\end{tabular}
\end{center}

\subsection{OpenStreetMap via OSMnx}

C'est la source centrale. OSMnx interroge l'API Nominatim (g\'{e}ocodage) et l'API Overpass
(r\'{e}seau routier) pour t\'{e}l\'{e}charger le graphe de la zone demand\'{e}e. Deux modes de
t\'{e}l\'{e}chargement sont impl\'{e}ment\'{e}s~:
\begin{itemize}
  \item par nom de lieu (ex. \og Le Marais, Paris \fg{})~;
  \item par point central + rayon en m\`{e}tres (mode pr\'{e}f\'{e}r\'{e} quand les coordonn\'{e}es
        GPS sont connues via Mapbox).
\end{itemize}
En cas d'\'{e}chec du t\'{e}l\'{e}chargement par nom, un fallback automatique tente le t\'{e}l\'{e}chargement
par point. Le graphe est sauvegard\'{e} localement au format GraphML pour \'{e}viter de re-t\'{e}l\'{e}charger.

\subsection{API Overpass (noms des clients)}

Apr\`{e}s la s\'{e}lection des n\oe uds du graphe comme points de livraison, l'application enrichit
leurs noms en interrogeant l'API Overpass dans un rayon de 1\,500\,m autour du centre. Les noms
r\'{e}els d'\'{e}tablissements OSM (commerces, \'{e}quipements publics, sites touristiques) sont
utilis\'{e}s en priorit\'{e}. Trois endpoints Overpass sont test\'{e}s successivement en cas de
d\'{e}faillance.

\subsection{API Open-Meteo (m\'{e}t\'{e}o)}

La m\'{e}t\'{e}o r\'{e}elle est r\'{e}cup\'{e}r\'{e}e gratuitement et sans cl\'{e} via Open-Meteo.
Les codes WMO renvoy\'{e}s (0 = clair, 61--65 = pluie, 71--75 = neige, 95 = orage) sont convertis
en facteur multiplicatif appliqu\'{e} sur la matrice des temps de trajet~: $\times 1{,}0$ par
beau temps, $\times 1{,}3$ sous la pluie, $\times 2{,}0$ sous la neige, $\times 2{,}5$ par
temps d'orage.

\subsection{API ADEME ImpactCO2 (facteurs CO2)}

Lorsque l'option \og CO2 ADEME \fg{} est activ\'{e}e, le facteur d'\'{e}mission CO2 par
kilom\`{e}tre est r\'{e}cup\'{e}r\'{e} dynamiquement via l'API ADEME ImpactCO2
(\texttt{https://impactco2.fr/api/v1/transport}). L'utilisateur peut choisir entre voiture
thermique (120\,g/km), voiture \'{e}lectrique ou v\'{e}lo. Sans cette option, un facteur
local par d\'{e}faut est utilis\'{e} (v\'{e}hicule : 120\,g/km, v\'{e}lo : 8\,g/km, marche : 0).

\subsection{API Mapbox Geocoding (auto-compl\'{e}tion)}

L'auto-compl\'{e}tion de l'adresse dans le formulaire de l'interface utilise l'API
Mapbox Geocoding. Si le token Mapbox est absent, la saisie libre reste fonctionnelle
(la zone est transmise telle quelle \`{a} OSMnx).

\subsection{Donn\'{e}es de r\'{e}f\'{e}rence pr\'{e}-calcul\'{e}es}

Pour le jeu de donn\'{e}es initial (5\`{e}me arrondissement de Paris, 10 clients), la matrice
de temps a \'{e}t\'{e} g\'{e}n\'{e}r\'{e}e via l'\textbf{API publique OSRM}
\url{http://router.project-osrm.org/table/v1/driving/} et sauvegard\'{e}e dans
\texttt{core\_data/live\_time\_matrix.npy}. Pour toutes les missions dynamiques g\'{e}n\'{e}r\'{e}es
depuis l'interface, les matrices sont calcul\'{e}es \textbf{localement} \`{a} partir du graphe OSMnx,
sans appel OSRM.

%=======================================================
\section{Biblioth\`{e}ques, Outils et technologies}
%=======================================================

\subsection{Backend -- Python}

\begin{table}[h]
  \begin{center}
    \begin{tabular}{|l|l|}
      \hline
      \textbf{Biblioth\`{e}que} & \textbf{R\^{o}le} \\
      \hline
      \texttt{FastAPI} & Framework HTTP : routes REST et WebSockets \\
      \texttt{uvicorn[standard]} & Serveur ASGI \\
      \texttt{Pydantic v2} & Validation des entr\'{e}es (mod\`{e}les de requ\^{e}te/r\'{e}ponse) \\
      \texttt{ortools} & Solveur VRPTW de Google \\
      \texttt{osmnx} $\geq$2.0.0 & T\'{e}l\'{e}chargement et manipulation des graphes OSM \\
      \texttt{networkx} & Algorithmes de graphes (SCC, Louvain, centralit\'{e}, etc.) \\
      \texttt{numpy} & Matrices de co\^{u}t et calculs num\'{e}riques \\
      \texttt{pandas} & Lecture et \'{e}criture des CSV de livraisons \\
      \texttt{requests} & Appels HTTP (Open-Meteo, Overpass, ADEME) \\
      \texttt{shapely} & Enveloppes convexes pour les secteurs Louvain \\
      \texttt{sqlite3} & Base de donn\'{e}es relationnelle embarqu\'{e}e \\
      \texttt{pytest} + \texttt{httpx} & Tests unitaires et d'int\'{e}gration \\
      \hline
    \end{tabular}
  \end{center}
  \caption{Biblioth\`{e}ques Python (fichier \texttt{requirements.txt})}
  \label{tab:python-libs}
\end{table}

\subsection{Frontend -- TypeScript / Next.js}

\begin{table}[h]
  \begin{center}
    \begin{tabular}{|l|l|l|}
      \hline
      \textbf{Biblioth\`{e}que} & \textbf{Version} & \textbf{R\^{o}le} \\
      \hline
      \texttt{next} & \textasciicircum14.2 & Framework React SSR \\
      \texttt{react} & \textasciicircum18.3 & Rendu UI \\
      \texttt{typescript} & \textasciicircum5.6 & Typage statique \\
      \texttt{@tanstack/react-query} & \textasciicircum5.59 & Cache et synchro des requ\^{e}tes API \\
      \texttt{leaflet} + \texttt{react-leaflet} & 1.9/4.2 & Carte interactive (tourn\'{e}es, markers) \\
      \texttt{mapbox-gl} & \textasciicircum3.21 & Tuiles de fond de carte \\
      \texttt{recharts} & \textasciicircum3.8 & Graphiques de comparaison KPI \\
      \texttt{next-auth} & \textasciicircum4.24 & Authentification OAuth \\
      \texttt{@playwright/test} & \textasciicircum1.54 & Tests E2E \\
      \hline
    \end{tabular}
  \end{center}
  \caption{D\'{e}pendances frontend (\texttt{frontend/package.json})}
  \label{tab:frontend-libs}
\end{table}

%=======================================================
\section{Travail r\'{e}alis\'{e}}
%=======================================================

\subsection{Construction du graphe}

\subsubsection{Structure}

Le graphe routier est un \textbf{networkx.MultiDiGraph dirig\'{e}} charg\'{e} depuis un fichier
GraphML (format standard OSM). Chaque n\oe ud est une intersection du r\'{e}seau routier avec
ses coordonn\'{e}es GPS. Chaque ar\^{e}te est un tron\c{c}on de voie avec ses attributs physiques
(longueur, type de voie) enrichis par annotation~:

\begin{table}[h]
  \begin{center}
    \begin{tabular}{|l|l|l|}
      \hline
      \textbf{Attribut} & \textbf{Unit\'{e}} & \textbf{Calcul} \\
      \hline
      \texttt{travel\_time} & secondes & longueur / vitesse (selon type de voie) \\
      \texttt{co2\_g} & grammes & (longueur / 1000) $\times$ facteur CO2 du mode \\
      \texttt{risk\_score} & sans unit\'{e} & (longueur / 1000) $\times$ facteur risque de la voie \\
      \texttt{length} & m\`{e}tres & donn\'{e}e OSM ou haversine si absente \\
      \hline
    \end{tabular}
  \end{center}
  \caption{Attributs calcul\'{e}s sur les ar\^{e}tes (\texttt{scripts/generator\_engine.py})}
  \label{tab:edge-attrs}
\end{table}

Les facteurs de risque varient selon le type de voie~: motorway (1,8), residential (0,9),
cycleway (0,55), footway (0,45). Les modes de transport configurent vitesse et CO2~:
voiture (35\,km/h, 120\,g/km), v\'{e}lo (18\,km/h, 8\,g/km), marche (5\,km/h, 0\,g/km).

\subsubsection{Graphe de r\'{e}f\'{e}rence (paris5)}

Le jeu de donn\'{e}es pr\'{e}-calcul\'{e} couvre un quartier de Paris (\texttt{core\_data/paris5.graphml}).
Ses m\'{e}triques structurelles sont document\'{e}es dans \texttt{core\_data/graph\_analysis\_soutenance.json}~:

\begin{table}[h]
  \begin{center}
    \begin{tabular}{|l|c|}
      \hline
      N\oe uds & 124 \\
      Ar\^{e}tes & 210 \\
      Plus court chemin moyen & 13,485 \\
      Diam\`{e}tre & 35 \\
      Densit\'{e} & 0,01377 \\
      Clustering moyen & 0,0685 \\
      \hline
    \end{tabular}
  \end{center}
  \caption{M\'{e}triques du graphe paris5}
  \label{tab:graph-metrics}
\end{table}

\subsubsection{Les matrices de co\^{u}t}

\`{A} partir du graphe annot\'{e}, plusieurs matrices carr\'{e}es $n \times n$ sont construites
et sauvegard\'{e}es au format NumPy (\texttt{.npy}). $n$ est le nombre total de points,
d\'{e}p\^{o}t inclus~; l'indice~0 est toujours le d\'{e}p\^{o}t. Le jeu de r\'{e}f\'{e}rence
historique \texttt{core\_data/} contient directement les matrices temps et distance
($n = 11$, d\'{e}p\^{o}t + 10 clients). Les missions g\'{e}n\'{e}r\'{e}es depuis l'interface
produisent aussi les matrices CO2, risque et composite dans leur dossier de cache.

\paragraph{Matrice de temps $T$ --- \texttt{live\_time\_matrix.npy}.}

Chaque entr\'{e}e $T_{ij}$ est le \textbf{temps de trajet minimal en secondes} pour aller
du point $i$ au point $j$ en suivant le plus court chemin pond\'{e}r\'{e} par
\texttt{travel\_time} sur le graphe OSM annot\'{e}. Pour le jeu de r\'{e}f\'{e}rence,
$T_{ij}$ varie de 13\,s (Boulangerie $\rightarrow$ Garage, distants de 87\,m) \`{a} 412\,s
(D\'{e}p\^{o}t $\rightarrow$ H\^{o}tel de Ville, 2\,312\,m).

La matrice est \textbf{asym\'{e}trique}~: $T_{ij} \neq T_{ji}$ en g\'{e}n\'{e}ral, car le
graphe OSM est dirig\'{e} (sens uniques, vitesses diff\'{e}rentes selon le sens de circulation).
Par exemple, $T_{0 \rightarrow 1} = 412$\,s mais $T_{1 \rightarrow 0} = 130$\,s.

Le calcul de chaque temps d'ar\^{e}te (\texttt{\_annotate\_multimodal\_edges}) est~:
\[
  T^{\text{ar\^{e}te}} = \frac{\text{longueur (m)}}{\text{vitesse (m/s)}}
\]
La vitesse est~: la limite OSM (\texttt{maxspeed}) si disponible, sinon la vitesse par
d\'{e}faut du type de voie (50\,km/h pour \texttt{tertiary}, 30\,km/h pour
\texttt{residential}, etc.). La vitesse minimale est born\'{e}e \`{a} 0,8\,m/s pour
\'{e}viter les valeurs nulles. Sur le jeu de r\'{e}f\'{e}rence, les vitesses implicites
mesur\'{e}es varient entre 17 et 24\,km/h, coh\'{e}rent avec la congestion urbaine parisienne.

\paragraph{Matrice de distances $D$ --- \texttt{matrix\_5eme.npy}.}

Chaque entr\'{e}e $D_{ij}$ est la \textbf{distance en m\`{e}tres} du plus court chemin
pond\'{e}r\'{e} par \texttt{length}. La matrice est aussi asym\'{e}trique~: le chemin le
plus court en distance n'est pas n\'{e}cessairement le m\^{e}me que le chemin le plus court
en temps. Sur le jeu de r\'{e}f\'{e}rence~: $D_{ij} \in [87\,\text{m},\; 2\,312\,\text{m}]$.

Lorsqu'une ar\^{e}te OSM n'a pas d'attribut \texttt{length} valide, la longueur est
calcul\'{e}e par la formule \textbf{haversine} \`{a} partir des coordonn\'{e}es des
extr\'{e}mit\'{e}s.

\paragraph{Matrice CO2 $K$ --- \texttt{co2\_matrix.npy}.}

Chaque entr\'{e}e $K_{ij}$ est l'\textbf{\'{e}mission de CO2 en grammes} pour aller de
$i$ \`{a} $j$. Elle est calcul\'{e}e par Dijkstra avec le poids \texttt{co2\_g},
c'est-\`{a}-dire en cherchant le chemin le moins \'{e}metteur selon le mod\`{e}le CO2~:
\[
  K_{ij} = \sum_{\text{ar\^{e}tes}(i \to j)} \frac{\text{longueur\_ar\^{e}te (m)}}{1000}
            \times f_{\text{CO2}}
\]
avec $f_{\text{CO2}} = 120{,}0$\,g/km par d\'{e}faut pour la voiture. En mode monomodal
avec un facteur CO2 constant, cela revient souvent \`{a} minimiser la distance~:
\[
  K_{ij} = \frac{D_{ij}}{1000} \times f_{\text{CO2}}
\]
En mode multimodal, les chemins peuvent toutefois diff\'{e}rer, car le v\'{e}lo et la
marche ont des facteurs d'\'{e}mission plus faibles que la voiture.
Dans les missions dynamiques, cette matrice est sauvegard\'{e}e dans le dossier propre
\`{a} la mission. Si l'option ADEME est activ\'{e}e, le facteur $f_{\text{CO2}}$ est
r\'{e}cup\'{e}r\'{e} dynamiquement depuis l'API ImpactCO2 et peut diff\'{e}rer
sensiblement de la valeur par d\'{e}faut.

Cette matrice permet au solveur de tenir compte des \'{e}missions lorsque
\texttt{optimization\_target = "composite"}. L'objectif \texttt{"distance"} peut aussi
r\'{e}duire indirectement le CO2 en mode monomodal, car l'\'{e}mission est alors
proportionnelle \`{a} la distance. Elle sert \'{e}galement au module d'analyse CO2 pour
calculer le bilan environnemental de chaque tourn\'{e}e.

\paragraph{Matrice de risque $R$ --- \texttt{risk\_matrix.npy}.}

Chaque entr\'{e}e $R_{ij}$ est un \textbf{score de risque sans unit\'{e}} repr\'{e}sentant
la dangerosit\'{e} de l'itin\'{e}raire. Pour chaque ar\^{e}te, le risque est pond\'{e}r\'{e}
par la longueur et le type de voie~:
\[
  R^{\text{ar\^{e}te}} = \frac{\text{longueur (m)}}{1000} \times f_{\text{risk}}
  \times f_{\text{oneway}}
\]
avec $f_{\text{oneway}} = 1{,}08$ si la voie est \`{a} sens unique et le mode est v\'{e}lo
ou marche (d\'{e}passement interdit, risque accru). Les facteurs par type de voie vont de
1,8 (autoroute) \`{a} 0,45 (chemin pi\'{e}ton), en passant par 0,9 (rue r\'{e}sidentielle).
Cette matrice n'est pas utilis\'{e}e seule comme objectif~; elle est uniquement un
composant de la matrice composite.

\paragraph{Matrice composite $C$ --- \texttt{composite\_cost\_matrix.npy}.}

C'est la matrice utilis\'{e}e par OR-Tools lorsque \texttt{optimization\_target = "composite"}.
Elle combine les quatre pr\'{e}c\'{e}dentes en une seule mesure de co\^{u}t global.

Avant la combinaison, chaque matrice est \textbf{normalis\'{e}e par sa m\'{e}diane}
(fonction \texttt{\_robust\_scale})~:
\[
  \tilde{M}_{ij} = \frac{M_{ij}}{\text{m\'{e}d}(M^+)}
\]
o\`{u} $M^+$ d\'{e}signe les valeurs finies positives de $M$. La normalisation par la
m\'{e}diane est plus robuste que la normalisation par le maximum~: elle n'est pas
influenc\'{e}e par les rares paires tr\`{e}s \'{e}loign\'{e}es ou par des valeurs aberrantes.
La matrice composite est ensuite~:
\[
  C_{ij} = w_t \cdot \tilde{T}_{ij} + w_d \cdot \tilde{D}_{ij}
           + w_c \cdot \tilde{K}_{ij} + w_r \cdot \tilde{R}_{ij}
\]
avec $w_t = 0{,}55$, $w_d = 0{,}20$, $w_c = 0{,}15$, $w_r = 0{,}10$.
Les poids sont normalis\'{e}s pour sommer \`{a} 1.

Dans les missions dynamiques, la valeur composite guide OR-Tools vers les clients
qui sont \`{a} la fois rapides,
proches, peu polluants et peu risqu\'{e}s.

\paragraph{Matrice d'incidents $T_{\text{inc}}$ --- \texttt{live\_time\_matrix\_incident.npy}.}

Une sixi\`{e}me matrice est g\'{e}n\'{e}r\'{e}e dynamiquement lorsqu'un incident est
simul\'{e} (\texttt{simulate\_incident\_replan()}). Elle est identique \`{a} $T$ sauf
pour les ar\^{e}tes affect\'{e}es par l'incident, dont le temps est multipli\'{e} par le
facteur de p\'{e}nalit\'{e} (1,5 par d\'{e}faut, configurable). Les effets en cascade
peuvent \^{e}tre drastiques~: le tronçon Garage $\to$ Boulangerie (13\,s normalement)
monte \`{a} 378\,s (+2\,808\,\%) quand l'ar\^{e}te interposée est bloqu\'{e}e et que
le d\'{e}tour oblige \`{a} un long contournement. OR-Tools est relancé sur cette matrice
modifi\'{e}e pour recalculer les tourn\'{e}es.

\paragraph{Algorithme de construction (r\'{e}sum\'{e}).}

Pour chaque source $i$, quatre Dijkstra sont ex\'{e}cut\'{e}s via
\texttt{nx.single\_source\_dijkstra\_path\_length}~:
un par m\'{e}trique (\texttt{travel\_time}, \texttt{length}, \texttt{co2\_g},
\texttt{risk\_score}).
Chaque matrice est initialis\'{e}e \`{a} $10^9$ (paires inaccessibles) et la diagonale
\`{a} 0. Complexit\'{e} totale~: $O\bigl(n \times (V + E) \log V\bigr)$.

La matrice composite est construite apr\`{e}s les quatre autres, par normalisation
et combinaison lin\'{e}aire.

\paragraph{Mode multimodal.}

Quand \texttt{transport\_mode = "multimodal"}, les matrices sont calcul\'{e}es
\emph{ind\'{e}pendamment} pour chaque mode (drive, bike, walk) puis fusionn\'{e}es
par \textbf{minimum} pour chaque paire $(i, j)$~:
\[
  M^{\text{multi}}_{ij} = \min_{m \in \{\text{drive, bike, walk}\}} M^{(m)}_{ij}
\]
Si aller en v\'{e}lo de $i$ \`{a} $j$ est plus rapide qu'en voiture, c'est la valeur
v\'{e}lo qui est retenue dans $T$, et la valeur v\'{e}lo (8\,g/km) dans $K$.
Cela permet au solveur de proposer naturellement l'\'{e}co-mobilit\'{e} quand elle est
avantageuse.

\subsubsection{Ajustements post-construction}

\begin{itemize}
  \item \textbf{M\'{e}t\'{e}o}~: $T \leftarrow T \times \dfrac{f_{\text{m\'{e}t\'{e}o}}}{f_{\text{vitesse}}}$
        (ex. $\times 2$ sous la neige). En cas de conditions s\'{e}v\`{e}res,
        $\lfloor n/5 \rfloor$ indices al\'{e}atoires re\c{c}oivent une majoration
        suppl\'{e}mentaire $\times 1{,}5$ (perturbations locales).
  \item \textbf{Congestion horaire}~: $T \leftarrow T \times f_{\text{trafic}(h)}$
        ($\times 1{,}7$ \`{a} 8h, $\times 1{,}8$ \`{a} 18h, $\times 1{,}0$ en heures creuses).
  \item \textbf{Relief SRTM} (optionnel)~: les temps et distances sont ajust\'{e}s
        en fonction du d\'{e}nivel\'{e} calcul\'{e} par le module \texttt{elevation\_engine}.
\end{itemize}

\subsection{Solveur VRPTW (Google OR-Tools)}

\subsubsection{Principe de r\'{e}solution}

Google OR-Tools est un solveur de recherche op\'{e}rationnelle open-source. Pour le VRPTW,
il utilise un \textbf{RoutingModel} qui mod\'{e}lise le probl\`{e}me comme un graphe de d\'{e}cision~:
\begin{itemize}
  \item Chaque n\oe ud correspond \`{a} un client ou au d\'{e}p\^{o}t.
  \item Les co\^{u}ts de transit entre n\oe uds sont fournis par des \textit{callbacks} (fonctions
        Python qui lisent la matrice de co\^{u}t choisie).
  \item Les contraintes (capacit\'{e}, fenêtres de temps, horizon maximal) sont d\'{e}clar\'{e}es
        comme dimensions du mod\`{e}le.
  \item Un \textit{drop penalty} p\'{e}nalise les clients non servis~; OR-Tools peut en d\'{e}poser
        certains si la faisabilit\'{e} l'exige.
\end{itemize}

\subsubsection{Strat\'{e}gies disponibles}

Cinq strat\'{e}gies de premi\`{e}re solution sont configur\'{e}es~:
\begin{itemize}
  \item \textbf{PATH\_CHEAPEST\_ARC}~: construit la solution en ajoutant l'arc le moins cher.
  \item \textbf{PARALLEL\_CHEAPEST\_INSERTION}~: ins\`{e}re les clients en parall\`{e}le au
        co\^{u}t d'insertion minimal.
  \item \textbf{SAVINGS} (Clarke-Wright)~: fusionne des routes en partant de tourn\'{e}es
        aller-retour, selon les \'{e}conomies r\'{e}alis\'{e}es.
  \item \textbf{CHRISTOFIDES}~: construction inspir\'{e}e du TSP m\'{e}trique, utilis\'{e}e ici
        comme strat\'{e}gie initiale possible.
  \item \textbf{GLOBAL\_CHEAPEST\_ARC}~: variante globale de PATH\_CHEAPEST\_ARC.
\end{itemize}

Trois m\'{e}taheuristiques de recherche locale sont utilis\'{e}es pour am\'{e}liorer la solution
initiale~: Guided Local Search, Simulated Annealing, Tabu Search.

\subsubsection{Flotte dynamique et co\^{u}t fixe}

Le solveur ne fixe pas arbitrairement le nombre de tra\^{i}neaux. Il calcule d'abord
deux rep\`{e}res d\'{e}terministes~:
\[
  k_{\min} = \left\lceil \frac{\text{poids total}}{\text{capacit\'{e} d'un tra\^{i}neau}} \right\rceil
  \qquad
  k_{\text{base}} = \left\lceil \frac{n_{\text{clients}}}{3} \right\rceil
\]
$k_{\min}$ est le minimum physique impos\'{e} par la capacit\'{e}. $k_{\text{base}}$
est une heuristique prudente li\'{e}e au nombre de clients. Le backend teste ensuite
plusieurs valeurs de $k$ r\'{e}parties entre ces rep\`{e}res, par exemple
$142, 206, 270, 334$ pour une mission de 1000 clients.

Pour chaque valeur de $k$, le solveur lance une r\'{e}solution et calcule un co\^{u}t
global~:
\[
  \text{score}_{k} = \text{co\^{u}t op\'{e}rationnel} + \text{co\^{u}t non-livr\'{e}s}
\]
\begin{align*}
  \text{co\^{u}t op\'{e}rationnel}
    &= \text{temps} + \text{distance} + \text{co\^{u}t flotte} \\
  \text{co\^{u}t non-livr\'{e}s}
    &= \text{drop\_penalty} \times n_{\text{non livr\'{e}s}}
\end{align*}
Le nombre de tra\^{i}neaux retenu est donc celui qui offre le meilleur compromis
entre couverture des colis et rentabilit\'{e}.

\subsubsection{Limite de temps adaptative}

Toutes les missions ne requi\`{e}rent pas le m\^{e}me budget de calcul. La limite de temps
de r\'{e}solution s'adapte au nombre de clients~:
\[
  t_{\max} = \text{clamp}\!\left(8,\; t_{\text{base}} \times \sqrt{\frac{n}{20}},\; 45\right)
  \text{ secondes}
\]
Les petites missions sont r\'{e}solues en quelques secondes~; les grandes bénéficient de
plus de budget de recherche.

\subsubsection{Post-traitement}

Apr\`{e}s la r\'{e}solution OR-Tools, les tourn\'{e}es sont am\'{e}lior\'{e}es par des heuristiques
de recherche op\'{e}rationnelle (\texttt{scripts/ro\_improvements.py})~:
\textbf{ALNS} puis \textbf{ILS}. ALNS effectue des destructions/r\'{e}parations de
tourn\'{e}es (random, worst, related, greedy, regret-2). ILS applique ensuite une
recherche locale it\'{e}r\'{e}e avec des mouvements \textbf{3-opt}, \textbf{OR-opt},
\textbf{2-opt*} et \textbf{double-bridge}. Une v\'{e}rification d'invariants garantit
l'absence de doublons ou de clients manquants~; en cas de violation, la solution
OR-Tools brute est restaur\'{e}e.

\subsubsection{Recherche automatique du nombre de tra\^{i}neaux}

Un m\'{e}canisme de sonde (\textit{sleigh search}) lance plusieurs r\'{e}solutions en parall\`{e}le
avec diff\'{e}rentes valeurs de $k$ (nombre de v\'{e}hicules), puis s\'{e}lectionne le $k$ qui
minimise le co\^{u}t global d\'{e}crit ci-dessus. Ce r\'{e}sultat est affich\'{e} dans la
section \og S\'{e}lection auto des tra\^{i}neaux \fg{} de la page de r\'{e}sultats.

\subsubsection{Mode grande \'{e}chelle}

\`{A} partir d'un seuil de 150 clients, le projet n'utilise plus directement le solveur
classique OR-Tools sur l'instance compl\`{e}te. Cette limite \'{e}vite de faire exploser
le temps de calcul quand le nombre de combinaisons devient trop grand.

Le mode \textit{large scale} transforme le probl\`{e}me en deux \'{e}tapes~:
\begin{enumerate}
  \item \textbf{G\'{e}n\'{e}ration de tourn\'{e}es candidates}~: le solveur construit de nombreuses
        petites tourn\'{e}es avec une logique gloutonne locale. Depuis un client, il regarde
        les voisins les plus proches dans la matrice de temps, puis ajoute le meilleur client
        faisable.
  \item \textbf{S\'{e}lection des tourn\'{e}es}~: seules les candidates qui respectent le d\'{e}p\^{o}t,
        la capacit\'{e}, les fen\^{e}tres horaires et la limite de 8h sont conserv\'{e}es.
        Un mod\`{e}le CP-SAT choisit ensuite la meilleure combinaison de candidates, en
        \'{e}vitant qu'un client soit livr\'{e} deux fois et en p\'{e}nalisant les clients non
        livr\'{e}s.
\end{enumerate}

Si CP-SAT n'est pas disponible ou ne trouve pas de solution faisable, un fallback glouton
s\'{e}lectionne les routes qui couvrent le plus de nouveaux clients pour le meilleur ratio
couverture/co\^{u}t. Ce mode permet de traiter des missions de plusieurs centaines ou
milliers de colis tout en conservant les contraintes principales du VRP.

\subsection{Profils IA}

Les profils d\'{e}finissent la strat\'{e}gie de r\'{e}solution utilis\'{e}e par OR-Tools.
L'interface expose \textbf{trois profils} \`{a} l'utilisateur~; sept existent dans le backend.

\begin{table}[h]
  \begin{center}
    \begin{tabular}{|l|l|l|c|l|}
      \hline
      \textbf{Profil (UI)} & \textbf{Cible} & \textbf{Premi\`{e}re strat.} & \textbf{Bonus} & \textbf{M\'{e}ta} \\
      \hline
      Express & temps & parallel cheapest insertion & +2 & Guided local search \\
      \'{E}colo & distance & savings & +3 & Simulated annealing \\
      Prudent & temps + marges & parallel cheapest insertion & +4 & Guided local search \\
      \hline
    \end{tabular}
  \end{center}
  \caption{Profils IA expos\'{e}s dans l'interface (\texttt{frontend/app/solver/page.tsx})}
  \label{tab:profiles-ui}
\end{table}

Les profils suppl\'{e}mentaires du backend (Opportuniste, Agressive, Championne,
Championne Secteurs) permettent des strat\'{e}gies plus sp\'{e}cialis\'{e}es~: la Championne
dispose d'une limite de temps de 35 secondes et du bonus de difficult\'{e} le plus \'{e}lev\'{e}
(+8\,pts), tandis que Championne Secteurs (\texttt{spatial\_sectorization = True}) applique
en plus un K-Means spatial avant la r\'{e}solution.

\subsection{Algorithmes de graphes}

Il est essentiel de distinguer les algorithmes \textbf{utilis\'{e}s dans le pipeline de calcul}
(toute mission les ex\'{e}cute) de ceux utilis\'{e}s uniquement \`{a} des fins d'\textbf{analyse
ou de p\'{e}dagogie} (appels \`{a} la demande, sans impact sur la tourn\'{e}e calcul\'{e}e).

\subsubsection{Algorithmes du pipeline technique}

\paragraph{Dijkstra --- construction des matrices de co\^{u}t.}

\textbf{O\`{u}~:} fonction \texttt{\_compute\_matrices\_local()} dans \texttt{scripts/generator\_engine.py},
appel\'{e}e \`{a} chaque cr\'{e}ation de mission.

\textbf{Principe~:} pour chaque source $i$ parmi les $n$ points de livraison,
\begin{lstlisting}[language=Python]
nx.single_source_dijkstra_path_length(graph, i, weight=w)
\end{lstlisting}
est ex\'{e}cut\'{e} avec les quatre poids (\texttt{travel\_time}, \texttt{length},
\texttt{co2\_g}, \texttt{risk\_score}).
Cela produit les quatre matrices $n \times n$ n\'{e}cessaires au solveur.
Complexit\'{e} totale~: $O\bigl(n \times (V + E) \log V\bigr)$.

\textbf{Pourquoi Dijkstra et pas Floyd-Warshall~?}
Floyd-Warshall calcule \emph{tous} les plus courts chemins en $O(N^3)$ o\`{u} $N$ est
le nombre de n\oe uds du graphe routier complet (des milliers de n\oe uds pour une zone
urbaine). Dijkstra est ex\'{e}cut\'{e} seulement depuis les points utiles \`{a} la mission
(d\'{e}p\^{o}t et clients), ce qui est plus rapide en pratique que de calculer tous les couples
du graphe routier. Floyd-Warshall resterait adapt\'{e}
si l'on travaillait sur un graphe dense de petite taille~; ici le graphe OSM est creux
(densit\'{e} 0,01377) et la seule matrice utile est celle entre les $n$ points de livraison,
pas entre toutes les intersections.

\textbf{Pourquoi pas Bellman-Ford~?}
Tous les poids (temps, distance, CO2, risque) sont positifs ou nuls~; Dijkstra est donc
adapt\'{e} \`{a} ce cas. Bellman-Ford, conçu pour les poids n\'{e}gatifs, serait
$O(VE)$ -- plus lent pour rien.

\paragraph{A* avec heuristique Haversine --- reconstruction g\'{e}om\'{e}trique des segments.}

\textbf{O\`{u}~:} fonction \texttt{build\_ai\_payload()} dans \texttt{scripts/routing\_payloads.py},
appel\'{e}e \`{a} chaque r\'{e}solution pour afficher les tourn\'{e}es sur la carte.

OR-Tools retourne l'\textbf{ordre de visite} des clients (indices), mais pas le chemin routier
r\'{e}el entre eux. Pour tracer les trajets sur la carte Leaflet, A* reconstruit le chemin
r\'{e}el sur le graphe OSM pour chaque segment $(\text{client}_i \rightarrow \text{client}_{i+1})$~:
\[
  \texttt{route} = \texttt{nx.astar\_path}(\text{graph},\; \text{src},\; \text{dst},\;
  \text{heuristic}=h,\; \text{weight}=\texttt{"travel\_time"})
\]
L'heuristique Haversine~:
\[
  h(u, v) = \frac{\text{haversine}(u, v)}{13{,}89 \text{ m/s}}
\]
sert d'heuristique g\'{e}om\'{e}trique de guidage. Dans un contexte urbain born\'{e} autour
de 50\,km/h, elle se comporte comme une heuristique admissible pratique~; sur des voies
plus rapides, il faut plut\^{o}t la lire comme une approximation d'affichage. A* explore
souvent \textbf{moins de n\oe uds} gr\^{a}ce au guidage vers la destination.

\textbf{Pourquoi A* et pas Dijkstra pour la g\'{e}om\'{e}trie~?}
Les deux peuvent donner le m\^{e}me plus court chemin si le poids et l'heuristique sont
coh\'{e}rents.
A* est souvent plus rapide car l'heuristique Haversine lui permet d'\'{e}viter d'explorer
une partie des n\oe uds ``dans la mauvaise direction''.
Sur un graphe de quartier (124 n\oe uds pour paris5), le gain est modeste~; sur une
zone de plusieurs kilom\`{e}tres, il devient significatif. C'est aussi la raison pour
laquelle la page \textit{Explore} du site visualise la r\'{e}duction du nombre de n\oe uds
explor\'{e}s entre Dijkstra et A*.

Il faut toutefois distinguer son r\^{o}le de celui du solveur~: A* ne choisit pas les
tourn\'{e}es. Il sert \`{a} reconstruire ou proposer des chemins point-\`{a}-point sur la
carte. Les matrices d'entr\'{e}e du VRP sont, elles, construites principalement avec
Dijkstra.

\paragraph{A* et K plus courts chemins --- options de route pour le joueur humain.}

\textbf{O\`{u}~:} fonction \texttt{\_collect\_candidate\_routes()} dans \texttt{routing\_payloads.py}.

Quand le joueur choisit sa prochaine livraison, le backend propose jusqu'\`{a} 5 routes
alternatives. Ces candidates sont g\'{e}n\'{e}r\'{e}es en combinant~:
\begin{enumerate}
  \item \texttt{nx.astar\_path} sur le poids \texttt{travel\_time} (route la plus rapide)~;
  \item \texttt{ox.routing.shortest\_path} sur le poids \texttt{length} (route la plus courte)~;
  \item \texttt{ox.routing.k\_shortest\_paths} (algorithme de Yen) sur les deux poids,
        pour obtenir des alternatives structurellement diff\'{e}rentes.
\end{enumerate}
Un filtre de diversit\'{e} (\texttt{\_pick\_diverse\_options}) \'{e}limine ensuite les routes
trop similaires (ratio de chevauchement d'ar\^{e}tes $> 40\,\%$) pour que chaque option
pr\'{e}sent\'{e}e soit r\'{e}ellement diff\'{e}rente.

\textbf{Pourquoi l'algorithme de Yen~?}
Le joueur ne veut pas trois fois le m\^{e}me chemin l\'{e}g\`{e}rement modifi\'{e}~:
les $K$ plus courts chemins au sens de Yen garantissent que les routes sont
syst\'{e}matiquement diff\'{e}rentes en termes de co\^{u}t, offrant de vraies alternatives.

\paragraph{ALNS + ILS --- post-traitement des tourn\'{e}es OR-Tools.}

\textbf{O\`{u}~:} fonction \texttt{\_postprocess\_vrp\_tours()} dans
\texttt{final\_scripts/solve\_santa\_final.py}, ex\'{e}cut\'{e}e \emph{automatiquement}
apr\`{e}s chaque r\'{e}solution OR-Tools classique (uniquement pour les missions monomodales).

OR-Tools peut se retrouver bloqu\'{e} dans un optimum local si son budget de temps
est \'{e}puis\'{e}. Deux \'{e}tapes de post-traitement am\'{e}liorent la solution~:

\textbf{Phase 1 --- ALNS} (Ropke \& Pisinger, 2006)~:
\begin{itemize}
  \item 3 op\'{e}rateurs de \emph{destruction}~: suppression al\'{e}atoire, suppression
        des pires clients (``worst removal''), suppression par similarit\'{e} de Shaw.
  \item 2 op\'{e}rateurs de \emph{r\'{e}paration}~: r\'{e}insertion gloutonne, r\'{e}insertion
        par regret-2 (ins\`{e}re le client pour lequel l'\'{e}cart entre le meilleur et
        le deuxi\`{e}me meilleur emplacement est maximal).
  \item S\'{e}lection \textbf{adaptative par roulette}~: les op\'{e}rateurs qui ont
        r\'{e}cemment am\'{e}lior\'{e} la solution voient leur probabilit\'{e} de s\'{e}lection augmenter.
  \item Crit\`{e}re d'acceptation~: \textbf{recuit simul\'{e}} (temp\'{e}rature initiale = 5\,\%
        du co\^{u}t initial, refroidissement $\times 0{,}995$/it\'{e}ration).
  \item Nombre d'it\'{e}rations~: $\max(120, \min(400, 6n))$.
  \item Complexit\'{e} par it\'{e}ration~: $O(n^2)$.
\end{itemize}

\textbf{Phase 2 --- ILS} (Recherche locale it\'{e}r\'{e}e)~:
\begin{itemize}
  \item 20 it\'{e}rations de perturbation/am\'{e}lioration.
  \item Chaque it\'{e}ration~: perturbation \textbf{double-bridge} sur la route la plus
        longue, puis recherche locale compl\`{e}te (3-opt + OR-opt + 2-opt*).
  \item Double-bridge~: coupe la tourn\'{e}e en 4 segments et les r\'{e}assemble dans
        un ordre diff\'{e}rent --- garantit une diversification qu'aucun 2-opt
        ou 3-opt ne peut d\'{e}faire directement.
  \item Crit\`{e}re d'acceptation~: am\'{e}lioration stricte ($\Delta < -1$\,s).
\end{itemize}

\textbf{Pourquoi ALNS + ILS apr\`{e}s OR-Tools~?}
OR-Tools fait une excellente exploration initiale mais travaille dans l'espace de
repr\'{e}sentation de son mod\`{e}le de routage. ALNS et ILS travaillent directement
sur la structure des tourn\'{e}es avec des mouvements plus larges (\'{e}change de
clients entre v\'{e}hicules, double-bridge), ce qui permet d'\'{e}chapper aux optima
locaux que les m\'{e}taheuristiques internes d'OR-Tools ne peuvent pas franchir dans
leur budget de temps.

\subsubsection{Algorithmes d'analyse (sur demande)}

Ces algorithmes sont ex\'{e}cut\'{e}s uniquement quand l'utilisateur consulte certaines
pages ou appelle certaines routes API~; ils n'influencent pas la tourn\'{e}e calcul\'{e}e.

\paragraph{Centralit\'{e} d'interm\'{e}diarit\'{e} --- identification des n\oe uds critiques.}

\textbf{O\`{u}~:} route \texttt{GET /api/missions/\{id\}/graph/metrics},
impl\'{e}ment\'{e}e dans \texttt{scripts/ro\_improvements.py} (\texttt{compute\_graph\_metrics}).

\[
  BC(v) = \sum_{s \neq v \neq t} \frac{\sigma_{st}(v)}{\sigma_{st}}
\]

Calcul approximatif par \'{e}chantillonnage ($k = \min(30, n)$ pivots) pour rester
temps-r\'{e}el. Identifie les ``goulots d'\'{e}tranglement'' du r\'{e}seau.

\textbf{Pourquoi cet algorithme~?}
Le solveur VRPTW ignore la topologie du graphe~; il travaille sur les matrices pr\'{e}-calcul\'{e}es.
La betweenness centrality r\'{e}pond \`{a} une autre question~: quels n\oe uds du r\'{e}seau
routier sont les plus vuln\'{e}rables~? Cette information est utile pour comprendre
les risques d'incidents.

\textbf{R\'{e}sultat v\'{e}rifi\'{e} sur paris5}~: le \textbf{Pont Marie} (n\oe ud 368142,
$BC = 0{,}484$) est le point le plus critique. Sa suppression r\'{e}duit la principale
SCC de 124 \`{a} 49 n\oe uds (--60,5\,\%), avec un rapport de d\'{e}vastation de 4,9
(attaque cibl\'{e}e 4,9$\times$ plus destructrice qu'une suppression al\'{e}atoire).

\paragraph{D\'{e}tection de communaut\'{e}s (Louvain) --- secteurs de livraison.}

\textbf{O\`{u}~:} route \texttt{GET /api/missions/\{id\}/graph/delivery-sectors},
impl\'{e}ment\'{e}e dans \texttt{scripts/community\_detection.py}.

L'algorithme de Louvain optimise la \textbf{modularit\'{e}} $Q$ du partitionnement~:
un bon partitionnement a plus d'ar\^{e}tes \textit{intra-cluster} que ne l'attendrait
un graphe al\'{e}atoire de m\^{e}me degr\'{e}. Appliqu\'{e} au graphe routier non-orient\'{e},
il d\'{e}couvre des quartiers naturels. Pour chaque cluster, une enveloppe convexe
(Shapely) est calcul\'{e}e pour l'affichage.

\textbf{Pourquoi Louvain plut\^{o}t que K-Means~?}
K-Means partitionne selon la proximit\'{e} g\'{e}ographique euclidienne~; Louvain respecte
la \textbf{connectivit\'{e} r\'{e}elle du r\'{e}seau routier}. Deux zones proches \`{a} vol
d'oiseau mais s\'{e}par\'{e}es par un mur ou un fleuve seront correctement mises dans des
clusters diff\'{e}rents par Louvain, pas par K-Means.

\paragraph{Floyd-Warshall, Dijkstra step-by-step, A* bidirectionnel, Dijkstra bidirectionnel --- visualisation p\'{e}dagogique.}

Ces quatre algorithmes sont expos\'{e}s via la page \textit{Explore} du site et des routes
API d\'{e}di\'{e}es. Ils ne sont pas appel\'{e}s lors du calcul d'une tourn\'{e}e.

\textbf{Floyd-Warshall} est expos\'{e} pour illustrer le calcul de \emph{tous} les plus
courts chemins simultan\'{e}ment~: il montre combien de paires $(i,j)$ seraient am\'{e}lior\'{e}es
en passant par un point interm\'{e}diaire. Dans le pipeline principal, cette \'{e}tape n'est pas
n\'{e}cessaire car les matrices sont d\'{e}j\`{a} construites par plus courts chemins sur le
graphe OSM complet.

\textbf{Dijkstra step-by-step}, \textbf{A* bidirectionnel} et \textbf{Dijkstra bidirectionnel}
permettent de comparer visuellement le nombre de n\oe uds explor\'{e}s selon la strat\'{e}gie
de recherche~: Dijkstra explore en cercles concentriques~; l'A* se concentre vers la cible~;
le bidirectionnel r\'{e}duit l'espace explor\'{e} en partant des deux extr\'{e}mit\'{e}s.

\paragraph{2-opt, OR-opt, Plus Proche Voisin --- analyse du bilan de mission.}

Ces trois algorithmes sont calcul\'{e}s dans \texttt{get\_debrief()} \textbf{uniquement
pour l'analyse post-mission}, pas pour l'optimisation~:
\begin{itemize}
  \item \textbf{2-opt et OR-opt}~: appliqu\'{e}s aux routes humaines pour calculer le
        \emph{potentiel d'am\'{e}lioration} que le joueur aurait pu r\'{e}aliser. Affich\'{e}
        dans la section \og Analyse \fg{} du bilan.
  \item \textbf{Nearest Neighbor (Plus Proche Voisin)}~: construit une tourn\'{e}e mono-v\'{e}hicule
        na\"ive TSP comme r\'{e}f\'{e}rence p\'{e}dagogique. Complexit\'{e} $O(n^2)$.
        Montre combien OR-Tools fait mieux qu'une heuristique gloutonne simple.
\end{itemize}

\subsection{Architecture modulaire}

Le projet est organis\'{e} en modules ind\'{e}pendants aux responsabilit\'{e}s clairement
d\'{e}limit\'{e}es. Tous gravitent autour d'un coeur central~: le \textbf{solveur VRP}.

\subsubsection{Le module central~: le solveur VRPTW}

\textbf{Fichier~:} \texttt{final\_scripts/solve\_santa\_final.py} (1\,437 lignes).

C'est le module principal du projet. Tous les autres modules existent pour lui fournir
des donn\'{e}es, pour am\'{e}liorer ses r\'{e}sultats, ou pour les pr\'{e}senter. Sans ce
solveur, aucune tourn\'{e}e optimis\'{e}e n'existe.

Le module s'articule en six phases~:
\begin{enumerate}
  \item \textbf{Chargement}~: lecture du CSV de points, chargement des matrices NumPy, application du facteur m\'{e}t\'{e}o.
  \item \textbf{Dispatch}~: en dessous du seuil grande \'{e}chelle, le solveur classique est utilis\'{e}; au-dessus de 150 clients, le mode \textit{large scale} prend le relais.
  \item \textbf{Mod\'{e}lisation OR-Tools}~: cr\'{e}ation du \texttt{RoutingIndexManager} et du \texttt{RoutingModel}, d\'{e}claration des callbacks de co\^{u}t, des contraintes de capacit\'{e} et de temps, du co\^{u}t fixe par v\'{e}hicule.
  \item \textbf{R\'{e}solution}~: strat\'{e}gie de premi\`{e}re solution + m\'{e}taheuristique, dans le budget de temps adaptatif.
  \item \textbf{Post-traitement}~: ALNS (Ropke \& Pisinger) puis ILS (double-bridge + 3-opt/OR-opt/2-opt*) pour affiner la solution.
  \item \textbf{Sauvegarde}~: s\'{e}rialisation JSON des tourn\'{e}es avec toutes les m\'{e}triques (temps, distance, CO2, v\'{e}hicules utilis\'{e}s).
\end{enumerate}

Le module est appel\'{e} exclusivement via la fonction \texttt{solve\_vrp()} depuis
\texttt{services.py}. Il est stateless~: il ne lit pas la base de donn\'{e}es et ne
d\'{e}pend que des fichiers de donn\'{e}es qui lui sont pass\'{e}s en param\`{e}tre.

\subsubsection{Module de g\'{e}n\'{e}ration de zone}

\textbf{Fichier~:} \texttt{scripts/generator\_engine.py} (989 lignes).

Ce module produit toutes les donn\'{e}es d'entr\'{e}e n\'{e}cessaires au solveur.
Il est d\'{e}clench\'{e} une seule fois par mission, \`{a} la cr\'{e}ation.

Son r\^{o}le~:
\begin{itemize}
  \item T\'{e}l\'{e}charger le graphe routier depuis OpenStreetMap (OSMnx).
  \item Annoter les ar\^{e}tes avec les quatre m\'{e}triques (temps, distance, CO2, risque).
  \item S\'{e}lectionner le d\'{e}p\^{o}t (noeud le plus proche du centre fourni par l'utilisateur) et les clients (noeuds al\'{e}atoires dans la zone).
  \item Enrichir les noms des clients via l'API Overpass.
  \item Calculer les quatre matrices de base par Dijkstra, puis construire la matrice composite.
  \item Appliquer les facteurs de congestion horaire et le relief SRTM si demand\'{e}.
  \item Sauvegarder graphe, CSV et matrices.
\end{itemize}

Sans ce module, il est impossible de g\'{e}n\'{e}rer une nouvelle zone~: le solveur
n'aurait pas de matrices \`{a} consommer.

\subsubsection{Module d'am\'{e}lioration des routes}

\textbf{Fichier~:} \texttt{scripts/ro\_improvements.py} (1\,803 lignes).

C'est le deuxi\`{e}me module le plus volumineux apr\`{e}s \texttt{services.py}. Il regroupe
les principaux algorithmes de graphes et d'optimisation combinatoire du projet.

\textbf{Fonctions utilis\'{e}es dans le pipeline de calcul~:}
\begin{itemize}
  \item \texttt{adaptive\_large\_neighborhood\_search()} --- ALNS post-solveur (phase~1 du post-traitement).
  \item \texttt{iterated\_local\_search()} --- ILS post-solveur (phase~2, avec double-bridge).
  \item \texttt{compute\_graph\_metrics()} --- m\'{e}triques structurelles (densit\'{e}, SCC, betweenness approxim\'{e}e).
\end{itemize}

\textbf{Fonctions utilis\'{e}es pour l'analyse post-mission (bilan)~:}
\begin{itemize}
  \item \texttt{two\_opt\_routes()}, \texttt{or\_opt\_routes()} --- mesure du potentiel non exploit\'{e} par le joueur humain.
  \item \texttt{nearest\_neighbor\_tour()} --- baseline TSP gloutonne.
  \item \texttt{optimality\_gap()} --- \'{e}cart \`{a} la borne inf\'{e}rieure th\'{e}orique.
\end{itemize}

\textbf{Fonctions utilis\'{e}es pour la p\'{e}dagogie (visualisation)~:}
\begin{itemize}
  \item \texttt{dijkstra\_steps()}, \texttt{bidirectional\_astar\_steps()}, \texttt{bidirectional\_dijkstra\_steps()} --- export pas-\`{a}-pas pour la page Explore.
  \item \texttt{floyd\_warshall()} --- calcul de tous les plus courts chemins sur la matrice de temps.
\end{itemize}

\subsubsection{Module de routage et d'\'{e}tat humain}

\textbf{Fichier~:} \texttt{scripts/routing\_payloads.py} (725 lignes).

Ce module g\`{e}re l'interface entre les d\'{e}cisions du joueur humain et le graphe OSM.

\textbf{Fonctions principales~:}
\begin{itemize}
  \item \texttt{\_collect\_candidate\_routes()} --- g\'{e}n\`{e}re le pool de routes candidates pour les options humaines (A* Haversine + K plus courts chemins).
  \item \texttt{compute\_route\_options()} --- filtre, classe et annote les routes candidates (plus rapide, plus courte, incident, diversit\'{e}).
  \item \texttt{build\_ai\_payload()} --- reconstruit la g\'{e}om\'{e}trie de chaque segment IA sur le graphe OSM via A* (pour l'affichage carte).
  \item \texttt{build\_human\_live\_stats()} --- calcule les statistiques en temps r\'{e}el de l'\'{e}tat humain (ETA, poids charg\'{e}, retards).
  \item \texttt{load\_graph()} --- charge un graphe GraphML et v\'{e}rifie/compl\`{e}te les attributs de vitesse.
\end{itemize}

Ce module maintient aussi deux caches LRU thread-safe~: un cache de graphes
(\texttt{GRAPH\_CACHE\_MAX = 32} entr\'{e}es) et un cache de routes candidates
(\texttt{ROUTE\_CANDIDATES\_CACHE\_MAX = 2\,048} entr\'{e}es), pour \'{e}viter de
recalculer les m\^{e}mes routes \`{a} chaque clic du joueur.

\subsubsection{Module de service (orchestrateur)}

\textbf{Fichier~:} \texttt{backend/app/services.py} (5\,921 lignes).

C'est le module le plus volumineux du projet. Il ne contient aucun algorithme de graphe
en propre~; son r\^{o}le est d'\textbf{orchestrer} tous les autres modules et d'exposer
une API coh\'{e}rente au controller FastAPI.

Il g\`{e}re notamment~:
\begin{itemize}
  \item La cr\'{e}ation de mission (appel \`{a} \texttt{generator\_engine}), la persistance en base et le cache fichier.
  \item La r\'{e}solution VRP~: s\'{e}lection de la strat\'{e}gie (\texttt{resolve\_ai\_strategy}), appel au solveur, calcul du benchmark, construction de la r\'{e}ponse.
  \item La recherche automatique du nombre de tra\^{i}neaux~: sondes parall\`{e}les sur plusieurs valeurs de $k$.
  \item Le mode humain~: gestion de l'\'{e}tat (segments valid\'{e}s, routes par tra\^{i}neau), calcul des options, simulation d'incidents.
  \item Les profils IA~: dictionnaire \texttt{AI\_PROFILE\_PRESETS} et fonction \texttt{resolve\_ai\_strategy()} qui adapte dynamiquement les param\`{e}tres selon le contexte.
  \item Le module d'apprentissage~: entra\^{i}nement, \'{e}valuation, recommandation de profil.
  \item Les fonctionnalit\'{e}s sociales~: amiti\'{e}s, messagerie directe, blocages.
  \item Le mode versus~: cr\'{e}ation de matchs, matchmaking, WebSockets.
\end{itemize}

\subsubsection{Module de persistance}

\textbf{Fichier~:} \texttt{backend/app/repository.py} (1\,844 lignes).

G\`{e}re la base de donn\'{e}es SQLite (\texttt{cache/api\_missions/operation\_noel.db}).
Toutes les requ\^{e}tes SQL sont isol\'{e}es dans ce module~; \texttt{services.py} ne
manipule jamais de SQL directement.

Tables principales~:
\begin{itemize}
  \item \texttt{players} --- profils, emails, hash de mot de passe (PBKDF2-HMAC-SHA256).
  \item \texttt{mission\_snapshots} --- cache des missions avec toutes leurs donn\'{e}es JSON (mission, r\'{e}sultats, benchmark, \'{e}tat humain, d\'{e}briefing).
  \item \texttt{leaderboard} --- scores tri\'{e}s par rang DESC.
  \item \texttt{versus\_matches}, \texttt{versus\_participants}, \texttt{versus\_invites}, \texttt{versus\_queue} --- syst\`{e}me de duels et matchmaking.
  \item \texttt{friendships}, \texttt{direct\_messages}, \texttt{user\_blocks} --- fonctionnalit\'{e}s sociales.
\end{itemize}

\subsubsection{Module API (contr\^{o}leur)}

\textbf{Fichier~:} \texttt{backend/app/main.py} (848 lignes).

D\'{e}finit les 70 routes HTTP REST et les 2 WebSockets. Son seul r\^{o}le est de~:
d\'{e}s\'{e}rialiser les requ\^{e}tes (Pydantic), appeler la fonction correspondante dans
\texttt{services.py}, et retourner la r\'{e}ponse ou une erreur HTTP structur\'{e}e.
Il ne contient aucune logique m\'{e}tier.

\subsubsection{Module de benchmark}

\textbf{Fichier~:} \texttt{scripts/benchmark\_engine.py} (148 lignes).

Calcule la solution na\"ive de r\'{e}f\'{e}rence (partition \'{e}gale des clients, visite
dans l'ordre d'index) et compare ses m\'{e}triques (temps, distance, CO2) avec la solution
OR-Tools. Produit le dictionnaire \texttt{savings} qui alimente la composante temps du score.

\subsubsection{Module m\'{e}t\'{e}o}

\textbf{Fichier~:} \texttt{scripts/weather\_engine.py} (83 lignes).

Deux modes~:
\begin{itemize}
  \item \texttt{get\_real\_weather()} --- g\'{e}ocode l'adresse via OSMnx puis interroge l'API Open-Meteo pour la m\'{e}t\'{e}o actuelle. Mappe les codes WMO vers les facteurs de la table \texttt{WEATHER\_SCENARIOS}.
  \item \texttt{get\_simulated\_weather()} --- tirage pond\'{e}r\'{e} al\'{e}atoire parmi les 7 sc\'{e}narios (plus de chances de beau temps).
\end{itemize}

\subsubsection{Module d'\'{e}l\'{e}vation SRTM (optionnel)}

\textbf{Fichier~:} \texttt{scripts/elevation\_engine.py} (208 lignes).

Interroge l'API OpenTopoData (donn\'{e}es SRTM NASA, r\'{e}solution 90\,m, gratuite) par
batches de 100 points. Calcule le d\'{e}nivel\'{e} entre chaque paire de points et ajuste
les matrices de temps et d'\'{e}nergie~:
\begin{itemize}
  \item Mont\'{e}e~: $\times (1 + \min(\text{pente} \times 6,\; 0{,}8))$ (jusqu'\`{a} +80\,\%)
  \item Descente~: $\times (1 + \max(\text{pente} \times 1{,}5,\; -0{,}2))$ (jusqu'\`{a} -20\,\%)
\end{itemize}
Activ\'{e} uniquement si l'utilisateur coche l'option \og Relief SRTM \fg{} dans le formulaire.

\subsubsection{Modules d'analyse du graphe}

\begin{itemize}
  \item \texttt{graph\_analysis.py} (110 lignes) --- calcule la centralit\'{e} de degr\'{e} et d'interm\'{e}diarit\'{e} (NetworkX). Produit le fichier \texttt{graph\_analysis\_soutenance.json}.
  \item \texttt{graph\_robustness.py} (138 lignes) --- mesure la d\'{e}gradation de la plus grande composante fortement connexe apr\`{e}s suppression cibl\'{e}e ou al\'{e}atoire de noeuds. Produit \texttt{robustness\_report.json}.
  \item \texttt{community\_detection.py} (103 lignes) --- Louvain sur le graphe non-orient\'{e}, enveloppes convexes Shapely.
  \item \texttt{zone\_clustering.py} (116 lignes) --- K-Means spatial impl\'{e}ment\'{e} de z\'{e}ro (sans scikit-learn).
\end{itemize}

Ces quatre modules sont invoqu\'{e}s soit par les routes API d'analyse, soit lors de
la g\'{e}n\'{e}ration du rapport technique automatis\'{e} (\texttt{generate\_technical\_report.py}).

\subsubsection{Modules utilitaires}

\begin{itemize}
  \item \texttt{mission\_paths.py} --- dataclass \texttt{MissionPaths} qui centralise tous les chemins de fichiers d'une mission (CSV, graphe, matrices, JSON). \'{E}vite les chemins ``en dur'' dispers\'{e}s dans le code.
  \item \texttt{osrm\_engine.py} --- script ponctuellement utilis\'{e} pour g\'{e}n\'{e}rer la matrice de temps OSRM du jeu de donn\'{e}es de r\'{e}f\'{e}rence (\texttt{live\_time\_matrix.npy}). Non utilis\'{e} pour les missions dynamiques.
  \item \texttt{main\_visualizer.py} --- g\'{e}n\'{e}rateur de carte Folium statique (HTML). Pr\'{e}c\'{e}dait la carte Leaflet interactive du frontend~; conserv\'{e} pour la visualisation hors-web.
\end{itemize}

\begin{table}[h]
  \begin{center}
    \footnotesize
    \begin{tabular}{|l|c|l|}
      \hline
      \textbf{Module} & \textbf{Lignes} & \textbf{R\^{o}le} \\
      \hline
      \texttt{services.py} & 5\,921 & Orchestrateur g\'{e}n\'{e}ral \\
      \textbf{\texttt{solve\_santa\_final.py}} & \textbf{1\,437} & \textbf{Solveur VRP (coeur)} \\
      \texttt{repository.py} & 1\,844 & Persistance SQLite \\
      \texttt{ro\_improvements.py} & 1\,803 & Algorithmes graphes + ALNS/ILS \\
      \texttt{generator\_engine.py} & 989 & G\'{e}n\'{e}ration de zone OSM \\
      \texttt{main.py} & 848 & Contr\^{o}leur FastAPI \\
      \texttt{routing\_payloads.py} & 725 & Routage humain + g\'{e}om\'{e}trie IA \\
      \texttt{schemas.py} & 338 & Mod\`{e}les Pydantic \\
      \texttt{elevation\_engine.py} & 208 & Relief SRTM \\
      \texttt{benchmark\_engine.py} & 148 & R\'{e}f\'{e}rence na\"ive \\
      \texttt{graph\_robustness.py} & 138 & Analyse de robustesse \\
      \texttt{community\_detection.py} & 103 & Secteurs Louvain \\
      \texttt{graph\_analysis.py} & 110 & Centralit\'{e} \\
      \texttt{zone\_clustering.py} & 116 & K-Means spatial \\
      \texttt{weather\_engine.py} & 83 & M\'{e}t\'{e}o (r\'{e}elle + simul\'{e}e) \\
      \texttt{mission\_paths.py} & 67 & Chemins de fichiers \\
      \hline
    \end{tabular}
  \end{center}
  \caption{Modules Python du projet, par taille}
  \label{tab:modules}
\end{table}

\subsection{Score et formules math\'{e}matiques}

\subsubsection{Solution de r\'{e}f\'{e}rence (na\"ive)}

Avant de calculer le score, il faut une r\'{e}f\'{e}rence. La \textbf{solution na\"ive} est
construite dans \texttt{scripts/benchmark\_engine.py} de la fa\c{c}on suivante~:
\begin{enumerate}
  \item Les $n$ clients sont r\'{e}partis en blocs \'{e}gaux entre les v\'{e}hicules avec
        \texttt{numpy.array\_split(clients, num\_vehicles)}.
  \item Dans chaque bloc, les clients sont visit\'{e}s dans \textbf{l'ordre de leur indice},
        sans aucune optimisation de l'ordre de visite.
  \item Le temps et la distance de chaque tron\c{c}on sont lus directement depuis les
        matrices pr\'{e}-calcul\'{e}es.
\end{enumerate}

Cette solution artificielle est volontairement simple~: elle sert de r\'{e}f\'{e}rence
non optimis\'{e}e pour mesurer le gain apport\'{e} par OR-Tools. Ce n'est pas une borne
inf\'{e}rieure math\'{e}matique, mais une baseline reproductible.

Les m\'{e}triques de la solution na\"ive sont calcul\'{e}es sur les m\^{e}mes matrices que
la solution optimis\'{e}e, pour une comparaison homog\`{e}ne. Les \'{e}conomies sont~:
\[
  \text{time\_saved\_pct} = \frac{T_{\text{na\"if}} - T_{\text{opt}}}{T_{\text{na\"if}}} \times 100
  \qquad \text{(born\'{e} \`{a} [-100, 100])}
\]
\[
  \text{co2\_saved\_kg} = \frac{K_{\text{na\"if}} - K_{\text{opt}}}{1000}
\]

\subsubsection{Calcul du score --- \'{e}tapes d\'{e}taill\'{e}es}

La fonction \texttt{computeScore()} du frontend et \texttt{get\_debrief()} du backend
utilisent la m\^{e}me formule, v\'{e}rifi\'{e}e coh\'{e}rente entre les deux.

\paragraph{\'Etape 1 --- Normalisation du gain temps.}
Sur les essais de r\'{e}f\'{e}rence, OR-Tools \'{e}conomise souvent du temps par rapport
\`{a} la solution na\"ive.
Pour que ce gain se traduise en score maximal de 45\,pts (composante temps~: 45\,\% du total),
on applique un facteur de normalisation de 2,5~:
\[
  S_t = \min\!\left(100,\; \text{time\_saved\_pct} \times 2{,}5\right)
\]
Ainsi, un gain de 40\,\% donne $S_t = 100$, un gain de 20\,\% donne $S_t = 50$.

\paragraph{\'Etape 2 --- Score CO2.}
Le CO2 \'{e}conom\'{e} est normalis\'{e} par une r\'{e}f\'{e}rence proportionnelle au nombre
de clients (0,1\,kg par client, minimum 1\,kg)~:
\[
  \text{co2\_ref} = \max(1{,}0,\; n \times 0{,}1)
  \qquad
  S_c = \min\!\left(100,\; \frac{\Delta\text{CO2 (kg)}}{\text{co2\_ref}} \times 100\right)
\]
Ce choix de r\'{e}f\'{e}rence \'{e}vite que les petites missions (peu de CO2 en valeur absolue)
soient syst\'{e}matiquement p\'{e}nalis\'{e}es.

\paragraph{\'Etape 3 --- Autres composantes.}
\begin{itemize}
  \item $S_b$ (budget)~: pourcentage du budget restant apr\`{e}s d\'{e}duction du co\^{u}t
        des tra\^{i}neaux utilis\'{e}s, born\'{e} dans $[0, 100]$.
  \item $S_v$ (couverture)~: $\dfrac{n_{\text{servis}}}{n_{\text{total}}} \times 100$.
        Les clients non servis (\textit{dropped}) p\'{e}nalisent directement ce score.
\end{itemize}

\paragraph{\'Etape 4 --- Score de base pond\'{e}r\'{e}.}
\[
  S_{\text{base}} = \underbrace{0{,}45}_{45\,\text{pts max}} \cdot S_t
                  + \underbrace{0{,}20}_{20\,\text{pts max}} \cdot S_c
                  + \underbrace{0{,}10}_{10\,\text{pts max}} \cdot S_b
                  + \underbrace{0{,}25}_{25\,\text{pts max}} \cdot S_v
\]
Le total th\'{e}orique est 100\,pts. Le choix des poids privil\'{e}gie le gain de temps
(45\,\%) et la couverture des livraisons (25\,\%), puis le bilan CO2 (20\,\%) et
le budget (10\,\%).

\paragraph{\'Etape 5 --- Bonus.}
Des points suppl\'{e}mentaires r\'{e}compensent la difficult\'{e} et les performances
sp\'{e}ciales~:
\begin{table}[h]
  \begin{center}
    \begin{tabular}{|p{0.24\textwidth}|p{0.48\textwidth}|p{0.18\textwidth}|}
      \hline
      \textbf{Bonus} & \textbf{Condition} & \textbf{Valeur} \\
      \hline
      Profil IA & selon le profil choisi & +2 (Express) \`{a} +10 (Championne Secteurs) \\
      Incidents & \texttt{random\_incidents = True} & +10 \\
      Humain bat l'IA & temps humain $<$ temps IA & +5 \\
      M\'{e}t\'{e}o difficile & $B_\text{m\'{e}t\'{e}o} = \max(0, (\text{factor}-1)\times 8)$ & 0 (clair) \`{a} +12 (tempête) \\
      \hline
    \end{tabular}
  \end{center}
  \caption{Bonus du score final}
  \label{tab:bonuses}
\end{table}

Le bonus m\'{e}t\'{e}o d\'{e}taille~: par temps clair (factor = 1,0), bonus = 0~; sous la
pluie (factor = 1,3), bonus = 2,4~; sous la neige (factor = 2,0), bonus = 8~; par orage
(factor = 2,5), bonus = 12.

\paragraph{\'Etape 6 --- Score final.}
\[
  S_{\text{final}} = \text{clamp}\!\left(0,\;
    S_{\text{base}} + B_{\text{profil}} + B_{\text{incidents}} + B_{\text{humain}} + B_{\text{m\'{e}t\'{e}o}},\;
    100\right)
\]

\paragraph{Exemple concret.}
Sur le jeu de r\'{e}f\'{e}rence (10 clients, profil Express, temps clair)~:
\begin{itemize}
  \item $\text{time\_saved\_pct} = 53{,}2\,\%$ $\Rightarrow$ $S_t = \min(100, 53{,}2 \times 2{,}5) = 100$
  \item $\Delta\text{CO2} = 0{,}95$\,kg, $\text{co2\_ref} = 1{,}0$ $\Rightarrow$ $S_c = 95$
  \item $S_{\text{base}} = 0{,}45 \times 100 + 0{,}20 \times 95 + 0{,}10 \times S_b + 0{,}25 \times 100$
  \item $B_{\text{profil}} = +2$ (Express), m\'{e}t\'{e}o bonus = 0
\end{itemize}

\begin{table}[h]
  \begin{center}
    \begin{tabular}{|c|c|c|}
      \hline
      \textbf{Score} & \textbf{Rang} & \textbf{Titre} \\
      \hline
      $\geq 85$ & S & Eco-Livreur L\'{e}gendaire \\
      $\geq 70$ & A & Chef Logisticien \\
      $\geq 50$ & B & Livreur Efficace \\
      $\geq 30$ & C & Apprenti P\`{e}re No\"{e}l \\
      $< 30$ & D & En formation \\
      \hline
    \end{tabular}
  \end{center}
  \caption{Syst\`{e}me de rangs}
  \label{tab:ranks}
\end{table}

\subsection{Recommandation contextuelle (AI Learning v2.0)}

\subsubsection{Principe et contrainte de donn\'{e}es}

Le syst\`{e}me d'apprentissage tente de r\'{e}pondre \`{a} la question~: \emph{quel profil IA
devrait produire un bon r\'{e}sultat pour une mission au contexte donn\'{e}~?} Il apprend sur
l'historique des missions r\'{e}solues sauvegard\'{e}es en base SQLite.

Il ne s'agit pas d'un mod\`{e}le de machine learning lourd de type r\'{e}gression,
for\^{e}t al\'{e}atoire ou r\'{e}seau de neurones. Ce choix est volontaire~: un tel mod\`{e}le
aurait besoin de beaucoup de missions r\'{e}elles, vari\'{e}es et correctement annot\'{e}es
pour produire des recommandations fiables. Le projet utilise donc plut\^{o}t une
recommandation statistique contextuelle, plus simple \`{a} expliquer et plus adapt\'{e}e
\`{a} un historique encore limit\'{e}.

\textbf{Contrainte importante}~: le mod\`{e}le ne peut s'entra\^{i}ner que si au moins
8 missions ont \'{e}t\'{e} r\'{e}solues et stock\'{e}es. C'est une limitation r\'{e}aliste~:
un mod\`{e}le entra\^{i}n\'{e} sur trop peu de donn\'{e}es ne serait pas fiable. En pratique,
plus l'historique est riche (missions dans des contextes vari\'{e}s~: pluie, neige, incidents,
villes diff\'{e}rentes), plus les recommandations sont pr\'{e}cises.

\subsubsection{Repr\'{e}sentation du contexte}

Chaque mission est r\'{e}sum\'{e}e par une cl\'{e} de contexte combinant~:
m\'{e}t\'{e}o (bucket), pr\'{e}sence d'incidents, taille de la mission (small/medium/large),
niveau du budget, co\^{u}t du tra\^{i}neau, densit\'{e} de clients. Cela permet de regrouper
les missions similaires et d'en tirer des leçons.

\subsubsection{Estimation bay\'{e}sienne avec lissage de Laplace}

Pour \'{e}viter de tout parier sur un contexte avec peu de donn\'{e}es, un lissage de Laplace
($\alpha = 3{,}0$) m\'{e}lange la moyenne locale (contexte sp\'{e}cifique) avec la moyenne
globale (tous contextes confondus)~:
\[
  \hat{\mu} = \frac{n_c \cdot \mu_c + \alpha \cdot \mu_g}{n_c + \alpha}
\]
Quand $n_c = 0$ (aucune donn\'{e}e dans ce contexte), le syst\`{e}me recourt enti\`{e}rement \`{a}
la moyenne globale. Quand $n_c$ cro\^{i}t, la moyenne locale prend progressivement le dessus.

\subsubsection{Score de confiance}

La recommandation est accompagn\'{e}e d'un score de confiance (0 \`{a} 0,95) qui tient compte
du nombre d'observations et de l'\'{e}cart entre le meilleur profil recommand\'{e} et le second~;
plus la marge est grande, plus la recommandation est fiable.

\subsubsection{Tuner OR-Tools}

Un second mod\`{e}le estime quels param\`{e}tres internes d'OR-Tools (strat\'{e}gie de premi\`{e}re
solution, m\'{e}taheuristique, limites de temps, etc.) semblent les plus performants pour un contexte donn\'{e}.
Il n\'{e}cessite au minimum 12 observations.

\subsection{Benchmark et comparaison}

La solution OR-Tools est syst\'{e}matiquement compar\'{e}e \`{a} une solution \textbf{na\"ive}~:
les clients sont r\'{e}partis en blocs \'{e}gaux entre les tra\^{i}neaux, dans l'ordre de
leur identifiant, sans aucune optimisation. Cette solution na\"ive sert de r\'{e}f\'{e}rence
pour calculer le gain. Les valeurs ci-dessous doivent \^{e}tre interpr\'{e}t\'{e}es comme un
exemple sur un jeu de r\'{e}f\'{e}rence donn\'{e}, et non comme une garantie de performance
g\'{e}n\'{e}rale sur toutes les villes ou toutes les tailles de mission.

\begin{table}[h]
  \begin{center}
    \begin{tabular}{|l|c|c|}
      \hline
      & \textbf{Solution na\"ive} & \textbf{OR-Tools} \\
      \hline
      Temps total & 2\,660\,s & 1\,245\,s \\
      Distance & 15\,293\,m & 7\,413\,m \\
      Gain temps & \multicolumn{2}{c|}{53,2\,\% (23 minutes)} \\
      CO2 \'{e}conom\'{e} & \multicolumn{2}{c|}{0,95 kg} \\
      \hline
    \end{tabular}
  \end{center}
  \caption{Benchmark sur le jeu de r\'{e}f\'{e}rence (10 clients, \texttt{core\_data/benchmark\_results.json})}
  \label{tab:benchmark}
\end{table}

%=======================================================
\section{Tests et validation}
%=======================================================

Le projet contient une suite de tests automatis\'{e}s qui vise surtout les parties \`{a} risque~:
solveur, contraintes, score, cache de graphes, API et parcours utilisateur.

\begin{itemize}
  \item \textbf{Tests backend Pytest}~: 20 fichiers couvrent les profils IA, le score,
        la s\'{e}lection du nombre de tra\^{i}neaux, le solveur grande \'{e}chelle,
        l'int\'{e}grit\'{e} du post-traitement, les routes API, l'authentification,
        la persistance SQLite et les options de route.
  \item \textbf{Tests frontend Playwright}~: 7 sc\'{e}narios E2E v\'{e}rifient les parcours
        de mission, les options de route, le routage libre, les \'{e}tats vides et une partie
        des parcours versus exp\'{e}rimentaux.
  \item \textbf{Tests cibl\'{e}s du solveur}~: le mode grande \'{e}chelle est test\'{e} avec
        des invariants simples~: d\'{e}part et retour au d\'{e}p\^{o}t, pas de client servi
        deux fois, respect de la capacit\'{e}, de la dur\'{e}e maximale et gestion possible
        des clients non livr\'{e}s.
\end{itemize}

Ces tests ne prouvent pas l'optimalit\'{e} globale des tourn\'{e}es, ce qui serait difficile
sur un probl\`{e}me NP-difficile. Ils servent plut\^{o}t \`{a} v\'{e}rifier que les contraintes
principales sont respect\'{e}es et que les r\'{e}gressions les plus probables sont d\'{e}tect\'{e}es.

%=======================================================
\section{Difficult\'{e}s rencontr\'{e}es}
%=======================================================

\subsection{Disponibilit\'{e} des APIs Open Data}

Les APIs publiques (Overpass, Open-Meteo) sont gratuites mais pas garanties. L'API Overpass
peut retourner des erreurs 429 (trop de requ\^{e}tes) ou \^{e}tre lente. Des m\'{e}canismes
de fallback ont \'{e}t\'{e} impl\'{e}ment\'{e}s~: basculement automatique entre trois endpoints
Overpass, et repli sur un graphe par point central si le t\'{e}l\'{e}chargement par nom de lieu
\'{e}choue.

\subsection{Faisabilit\'{e} du VRPTW}

Le VRPTW peut \^{e}tre infaisable si les contraintes sont trop strictes (trop peu de v\'{e}hicules
pour le nombre de clients, budget de temps insuffisant). La solution adopt\'{e}e est d'utiliser
un \textit{drop penalty}~: OR-Tools peut d\'{e}poser des clients plut\^{o}t que de rendre la
mission infaisable. Le score p\'{e}nalise les clients non servis via la composante de couverture.

\subsection{Contraintes arithm\'{e}tiques d'OR-Tools}

La p\'{e}nalit\'{e} de non-service doit rester dans les limites int64. Pour les grandes missions,
une formule dynamique garantit cette borne~: $\text{penalty} = \max(\text{base}, \text{horizon} \times n \times V)$,
born\'{e}e \`{a} $2 \times 10^9$.

\subsection{Apprentissage avec peu de donn\'{e}es}

Le module AI Learning n\'{e}cessite 8 missions minimum pour s'entra\^{i}ner. Les premi\`{e}res
utilisations de la plateforme ne b\'{e}n\'{e}ficient donc pas des recommandations personnalis\'{e}es.
Le lissage de Laplace ($\alpha = 3$) att\'{e}nue ce probl\`{e}me ``cold start'' en exploitant
la moyenne globale en l'absence de donn\'{e}es locales.

\subsection{Approximation des conditions r\'{e}elles}

Le projet cherche \`{a} se rapprocher du terrain, mais certaines donn\'{e}es restent
approximatives. La m\'{e}t\'{e}o est transform\'{e}e en facteurs multiplicatifs sur les temps
de trajet, les incidents sont simul\'{e}s, et le risque routier est estim\'{e} \`{a} partir
du type de voie. Ces choix permettent de tester plusieurs sc\'{e}narios sans d\'{e}pendre
de donn\'{e}es temps r\'{e}el complexes, mais ils ne remplacent pas une mesure directe de
la circulation ou des perturbations locales.

Une am\'{e}lioration importante serait donc de mieux calibrer ces facteurs~: comparer les
temps pr\'{e}dits aux temps observ\'{e}s, apprendre l'impact moyen de la pluie, de la neige
ou d'un incident sur plusieurs missions, puis ajuster automatiquement les coefficients.

%=======================================================
\section{Bilan}
%=======================================================

\subsection{Conclusion}

Op\'{e}ration No\"{e}l est une plateforme qui couvre le pipeline principal~:
depuis la saisie d'une adresse jusqu'\`{a} l'affichage d'une tourn\'{e}e optimis\'{e}e
sur une carte du monde r\'{e}el.

Sur le jeu de r\'{e}f\'{e}rence, les r\'{e}sultats sont encourageants~: OR-Tools r\'{e}alise
\textbf{53,2\,\%} d'\'{e}conomie de temps et \textbf{0,95\,kg} de CO2 \'{e}conom\'{e} par rapport
\`{a} la solution na\"ive. Ces chiffres restent li\'{e}s au sc\'{e}nario test\'{e}~; ils montrent
que l'approche fonctionne, mais ne constituent pas une garantie universelle. La plateforme
expose \'{e}galement plusieurs algorithmes de graphes classiques, visualisables interactivement.

Le projet r\'{e}pond donc \`{a} la probl\'{e}matique en montrant que les graphes et l'open data
permettent de construire une aide \`{a} la d\'{e}cision r\'{e}aliste pour les tourn\'{e}es du
P\`{e}re No\"{e}l. La limite principale reste la qualit\'{e} des approximations~: m\'{e}t\'{e}o,
incidents, vitesses et facteurs de risque doivent \^{e}tre mieux calibr\'{e}s pour rapprocher
encore le mod\`{e}le du terrain.

\subsection{Perspectives}

\begin{enumerate}
  \item \textbf{Fen\^{e}tres de temps plus r\'{e}alistes}~: les colonnes \texttt{tw\_start}
        et \texttt{tw\_end} sont d\'{e}j\`{a} exploit\'{e}es comme contraintes VRPTW.
        Une perspective serait de les rendre plus fines, par exemple selon le type
        de client, la priorit\'{e} du colis ou l'heure de d\'{e}part.
  \item \textbf{Accumulation d'historique}~: le mod\`{e}le AI Learning n\'{e}cessite un historique
        riche~; un m\'{e}canisme de pr\'{e}-chargement avec des missions synth\'{e}tiques dans des
        contextes vari\'{e}s am\'{e}liorerait les recommandations d\`{e}s la premi\`{e}re utilisation.
        Avec suffisamment de donn\'{e}es r\'{e}elles, il deviendrait pertinent de tester un
        mod\`{e}le de machine learning plus classique pour pr\'{e}dire le meilleur profil ou
        les meilleurs param\`{e}tres du solveur.
  \item \textbf{Calibration des approximations}~: la m\'{e}t\'{e}o, les incidents et le risque
        sont aujourd'hui mod\'{e}lis\'{e}s par des coefficients. Une perspective naturelle
        serait de les recalibrer \`{a} partir d'observations r\'{e}elles afin de rendre les
        temps de trajet et les scores plus robustes.
  \item \textbf{Passage \`{a} l'\'{e}chelle}~: le mode \textit{large scale} permet d\'{e}j\`{a}
        de traiter de grandes missions par g\'{e}n\'{e}ration de candidates et s\'{e}lection
        CP-SAT. La perspective principale est maintenant d'\'{e}valuer syst\'{e}matiquement
        la qualit\'{e} des solutions sur plusieurs villes et plusieurs densit\'{e}s.
  \item \textbf{V\'{e}hicules h\'{e}t\'{e}rog\`{e}nes}~: le solveur supporte d\'{e}j\`{a} un mode
        multimodal par v\'{e}hicule~; une interface de configuration explicite permettrait
        aux utilisateurs de mixer voiture, v\'{e}lo et marche dans la m\^{e}me mission.
  \item \textbf{Exposition des profils avanc\'{e}s}~: les profils Championne et Championne Secteurs
        ne sont pas accessibles depuis l'interface~; les exposer am\'{e}liorerait les scores
        obtenables et d\'{e}montrerait l'int\'{e}r\^{e}t du K-Means de pr\'{e}-sectorisation.
\end{enumerate}

%=======================================================
\section{Bibliographie}
\renewcommand{\bibname}{}
\renewcommand{\refname}{}
\begin{thebibliography}{9}

  \bibitem[COR09]{cormen}
    T.\,H. Cormen, C.\,E. Leiserson, R.\,L. Rivest, C. Stein,
    {\it Introduction to Algorithms}, 3\`{e}me \'{e}d., MIT Press, 2009.

  \bibitem[TOTH02]{toth}
    P. Toth, D. Vigo (eds.),
    {\it The Vehicle Routing Problem},
    SIAM, 2002.

  \bibitem[BOE17]{boeing}
    G. Boeing, ``OSMnx~: New Methods for Acquiring, Constructing, Analyzing, and
    Visualizing Complex Street Networks'',
    {\it Computers, Environment and Urban Systems}, 65, 2017.

  \bibitem[BLOND08]{louvain}
    V.\,D. Blondel et al.,
    ``Fast unfolding of communities in large networks'',
    {\it J. Stat. Mech.}, P10008, 2008.

\end{thebibliography}

%=======================================================
\section{Webographie}
\begin{thebibliography}{9}

  \bibitem[OSM]{osm}
    OpenStreetMap Foundation. \url{https://www.openstreetmap.org}

  \bibitem[OSMNXDOC]{osmnxdoc}
    G. Boeing, {\it OSMnx Documentation}. \url{https://osmnx.readthedocs.io}

  \bibitem[ORTOOLS]{ortools}
    Google Developers, {\it OR-Tools Vehicle Routing}.
    \url{https://developers.google.com/optimization/routing}

  \bibitem[OPENMETEO]{openmeteo}
    Open-Meteo, {\it Free Weather API}. \url{https://open-meteo.com}

  \bibitem[ADEME]{ademe}
    ADEME, {\it ImpactCO2 API}. \url{https://impactco2.fr/api/v1/transport}

  \bibitem[FASTAPI]{fastapi}
    S. Ram\'{i}rez, {\it FastAPI}. \url{https://fastapi.tiangolo.com}

  \bibitem[NEXTJS]{nextjs}
    Vercel, {\it Next.js 14}. \url{https://nextjs.org}

  \bibitem[NETWORKX]{networkx}
    A. Hagberg et al., {\it NetworkX}. \url{https://networkx.org}

  \bibitem[OVERPASS]{overpass}
    Overpass API. \url{https://overpass-api.de}

  \bibitem[MAPBOX]{mapbox}
    Mapbox, {\it Geocoding API}. \url{https://docs.mapbox.com/api/search/geocoding}

\end{thebibliography}

%=======================================================
\section{Annexes}
\appendix
\makeatletter
\def\@seccntformat#1{Annexe~\csname the#1\endcsname:\quad}
\makeatother

\section{Commandes utiles}

\begin{lstlisting}[language=bash]
# Installation et lancement
cd /home/bekkari/Documents/Graphes/Noel
docker-compose up --build

# Tests backend
PYTHONPATH=. pytest -q

# Tests frontend E2E
cd frontend
npx playwright test
\end{lstlisting}

\section{Parcours utilisateur synth\'{e}tique}

\begin{enumerate}
  \item Ouvrir \texttt{http://localhost:3000/solver}.
  \item Taper une adresse~: l'auto-compl\'{e}tion Mapbox propose des suggestions.
  \item R\'{e}gler le rayon, le nombre de colis, le budget, le profil IA.
  \item Lancer le calcul de la tourn\'{e}e~: la barre de progression
        affiche les \'{e}tapes (OSM, matrices, OR-Tools).
  \item Consulter la carte et le score. Simuler un incident si d\'{e}sir\'{e}.
  \item Acc\'{e}der \`{a} la page \texttt{Explore} pour visualiser Dijkstra, A* et Floyd-Warshall.
\end{enumerate}

\end{document}
